% =====================================================================
%  capitulos/cap10.tex — Capítulo 10. Las construcciones pronominales
%  Autoras: Mabel Giammatteo y Mariana Cuñarro
%  Fuente: Capitulo_10_normalizado.docx
% =====================================================================
\capitulo{Las construcciones pronominales}{Mabel Giammatteo y Mariana Cuñarro}

\section{Introducción}

En este capítulo nos ocupamos de los casos en que verbos como \forma{lavar,
divertir, abrazar} o \forma{acercar} pueden también emplearse en forma
pronominal ---\forma{lavarse, divertirse, abrazarse} o \forma{acercarse}---, lo
cual implica que el verbo y el pronombre deben coincidir en número y persona.
Para ello, en primer lugar, revisamos el funcionamiento de los pronombres
personales a fin de distinguir las características de aquellos que intervienen
en estas construcciones. Luego consideramos los dos grandes grupos en que se
dividen las construcciones pronominales: las \termino{reflejas}, en las que el
pronombre manifiesta un argumento del verbo, por lo que tiene papel temático y
desempeña la función sintáctica de COD o de COI, y las
\termino{cuasirreflejas}, en las que, aunque el pronombre comparte los rasgos
flexivos de número y persona del verbo, no representa un papel temático ni tiene
función sintáctica. No obstante, en todos los casos, la incorporación del
pronombre altera en diferentes sentidos el significado del verbo, ya sea porque
el evento que describe revierte sobre quien o quienes lo realizan (\forma{me
lavo, nos abrazamos}), o bien porque se producen cambios en la estructura
argumental (\forma{El viento voló los techos} frente a \forma{Los techos se
volaron}), la estructura eventiva (\forma{Comió una manzana} frente a \forma{Se
comió una manzana}) o la estructura sintáctica (\forma{Olvidé la clave} frente a
\forma{Me olvidé de la clave}).

\section{Los pronombres personales}\label{sec:10-pronombres}

Los pronombres personales son \termino{deícticos}: su referencia es variable de
acuerdo con los participantes del coloquio. «La primera persona se refiere al
hablante o a los hablantes, la segunda, al oyente o a los oyentes, y la tercera,
a las personas o cosas de las que se habla» \parencite[§\,4.4a]{rae2025}.

\begin{destacado}
En español, además de la categoría lingüística de \termino{persona} (primera,
segunda y tercera), los pronombres personales manifiestan \termino{número}
(singular y plural), \termino{género} (masculino y femenino) y \termino{caso}
(nominativo o caso del sujeto; acusativo o caso del OD; dativo o caso del OI, y
preposicional o terminal, que es el caso del pronombre cuando sigue a una
preposición).
\end{destacado}

\begin{table}[htbp]
\centering\footnotesize\setlength{\tabcolsep}{3pt}
\caption{Formas pronominales del español}\label{tab:10-pronombres}
\resizebox{\linewidth}{!}{%
\begin{tabular}{@{}l l c c c c c c c c@{}}
  \toprule
  & & \multicolumn{2}{c}{\textbf{Primera}} & \multicolumn{4}{c}{\textbf{Segunda}} &
    \multicolumn{2}{c}{\textbf{Tercera}} \\
  \cmidrule(lr){3-4}\cmidrule(lr){5-8}\cmidrule(l){9-10}
  & & & & \multicolumn{2}{c}{$-$cortesía\textsuperscript{a}} &
    \multicolumn{2}{c}{+cortesía} & & \\
  \cmidrule(lr){5-6}\cmidrule(lr){7-8}
  \textbf{Caso} & & Sg. & Pl. & Sg. & Pl. & Sg. & Pl. & Sg. & Pl. \\
  \midrule
  Nominativo & & yo & nosotros/-as & vos / tú & vosotros/-as\textsuperscript{b} &
    usted & ustedes & él / ella & ellos / ellas \\
  \addlinespace
  Objetivo\textsuperscript{c} & oblicuas: acusativo & \multirow{2}{*}{me} &
    \multirow{2}{*}{nos} & \multirow{2}{*}{te} & \multirow{2}{*}{os} & lo / la &
    los / las & lo / la & los / las \\
  & oblicuas: dativo & & & & & le & les & le & les \\
  & reflexiva & & & & & \multicolumn{2}{c}{se} & \multicolumn{2}{c}{se} \\
  \addlinespace
  Preposicional & oblicuas & mí & nosotros/-as & vos / ti & vosotros/-as & usted
    & ustedes & él / ella & ellos / ellas \\
  & reflexiva & & & & & \multicolumn{2}{c}{sí} & \multicolumn{2}{c}{sí} \\
  \bottomrule
\end{tabular}}

\medskip
\begin{minipage}{\linewidth}\scriptsize
\textsuperscript{a} Como señala \textcite[169]{ditullio2014}, además de las
categorías de persona, número, género y caso, hay una dimensión adicional,
presente en «los pronombres que designan al destinatario, la cortesía, que opone
las formas de tratamiento \textit{vos} o \textit{tú} (de acuerdo con el dialecto)
a \textit{usted} y, en español peninsular, \textit{vosotros} a \textit{ustedes}».
\textsuperscript{b} Las formas \textit{vosotros/-as} y \textit{os} se usan solo
en la variedad peninsular. \textit{Tú} se emplea en España y en muchas regiones
de Latinoamérica. \textit{Vos} es la forma típica de la variedad rioplatense,
pero también se utiliza en otras regiones de América, aunque con distinto grado
de aceptación social. \textsuperscript{c} Con la denominación de caso objetivo
reunimos al acusativo y al dativo, que solo tienen manifestación diferenciada en
la tercera persona, singular y plural, y también en las formas de segunda
persona \textit{usted} y \textit{ustedes}.
\end{minipage}
\end{table}

\subsection{Pronombres tónicos y pronombres átonos o clíticos}\label{sec:10-tonicos}

Los pronombres personales se pueden agrupar en dos series:

\begin{itemize}
  \item \termino{Tónicos}: llevan acento fuerte y son todos los que están en
        caso nominativo o preposicional.
  \item \termino{Átonos} o \termino{clíticos}: no tienen acento propio y se
        apoyan acentualmente en otra palabra a la que se anteponen
        (\termino{proclíticos}: \forma{Juan se cansó; ¿Me ayudás?}) o a la que
        se posponen (\termino{enclíticos}: \forma{Juan suele cansarse con
        facilidad; ¡Dame eso ya!}). Están en caso acusativo o dativo, y la
        primera y la segunda persona emplean la misma forma para ambos casos.
\end{itemize}

\begin{exe}
  \ex Después de dos horas de espera, el médico [\etq{Pr/COD} \textit{me}]
      atendió.
  \ex {}[\etq{Pr/COI} \textit{Me}] recetó unas pastillas para el dolor de
      cabeza.
\end{exe}

En la tercera persona, en cambio, las formas de acusativo y dativo difieren: se
emplean \forma{lo, la, los, las} para acusativo, y \forma{le} o \forma{les} para
dativo.

\begin{exe}
  \ex {}[\etq{Pr/COD} \textit{Lo}] encontré esta mañana en el bar de enfrente.
  \ex {}[\etq{Pr/COI} \textit{Le}] habían dado un premio en la Academia.
\end{exe}

\begin{ejercicio}\label{ej:10-caso}
En las siguientes oraciones señale el caso de los pronombres y justifique su
respuesta indicando la función sintáctica que cumplen, según se muestra en el
ejemplo resuelto~1.

\begin{exe}
  \ex \textit{Ud.} no sabe nada de lo sucedido.
  \ex ¿Cuántas veces \textit{lo} tengo que repetir?
  \ex Varios amigos \textit{me} comentaron del asunto.
  \ex Mis padres \textit{me} regalaron un coche nuevo.
  \ex \textit{Te lo} dije varias veces, pero no \textit{me} entendiste.
\end{exe}

\begin{resuelto}
1. \forma{Ud.}: caso nominativo, función sintáctica sujeto.
\end{resuelto}
\end{ejercicio}

\subsection{Clíticos oblicuos y reflexivos}\label{sec:10-oblicuos}

Cuando los clíticos difieren de la persona señalada por el sujeto ---o por la
flexión verbal, si el sujeto es tácito--- se dice que son \termino{oblicuos}
\exref{ex:10-medico}; en cambio, cuando son correferenciales con el sujeto, el
clítico es \termino{reflexivo} \exref{ex:10-yomisma}.

\begin{exe}
  \ex\label{ex:10-medico} {}[\etq{3.ª pers.} El médico] [\etq{1.ª pers.}
      \textit{me}] examinó las mamas cuidadosamente.
  \ex\label{ex:10-yomisma} {}[\etq{1.ª pers.} Yo misma] [\etq{1.ª pers.}
      \textit{me}] examiné las mamas cuidadosamente.
\end{exe}

Para que el pronombre sea reflexivo no solo debe coincidir con la persona del
sujeto, sino que ambos ---sujeto y objeto--- deben remitir al mismo referente
extralingüístico. En \exref{ex:10-paciente1}, aunque sujeto y objeto refieren a
la tercera persona, no hay reflexividad porque \forma{el médico} y \forma{la
paciente} son dos personas (es decir, referentes) distintas. Para indicar
correferencia en tercera persona se utiliza la forma reflexiva del pronombre:
\forma{se} \exref{ex:10-paciente2}.

\begin{exe}
  \ex\label{ex:10-paciente1} La paciente llegó al consultorio y [\etq{SN/Suj
      3.ª pers.} el médico] [\etq{Pr Oblic/COD 3.ª pers.} \textit{la}] examinó
      cuidadosamente.
  \ex\label{ex:10-paciente2} {}[\etq{SN/Suj 3.ª pers.} La paciente] [\etq{Pr
      Refl/COI 3.ª pers.} \textit{se}] examinó las mamas cuidadosamente.
\end{exe}

En el SV, los pronombres átonos representan a los argumentos internos del verbo;
por tanto, tienen papel temático en la estructura argumental (EA) y,
sintácticamente, cumplen las funciones de COD, \exref{ex:10-paciente1} y
\exref{ex:10-testigo}, o de COI, \exref{ex:10-paciente2} y
\exref{ex:10-vacuna}.

\begin{exe}
  \ex\label{ex:10-testigo} El testigo [\etq{Pr Refl/COD} \textit{se}] inculpó
      por el crimen.
  \ex\label{ex:10-vacuna} El investigador [\etq{Pr Refl/COI} \textit{se}]
      inyectó [\etq{SN/COD} la vacuna] en su propio brazo.
\end{exe}

\begin{ejercicio}\label{ej:10-funcion}
Determine la función sintáctica y el papel temático de los pronombres en
negrita, según se muestra en el ejemplo resuelto~1.

\begin{exe}
  \ex El jugo de naranja \textbf{le} desagrada.
  \ex Elena \textbf{le} dijo que esta vez llegaría temprano.
  \ex Elena \textbf{lo} dijo una vez más.
  \ex Elena \textbf{le} dijo a su hermano la verdad.
  \ex\label{ex:10-selo} Elena \textbf{se lo} dijo a su hermano.
\end{exe}

\begin{resuelto}
1. \forma{Le}: función sintáctica COI; papel temático experimentante.
\end{resuelto}
\end{ejercicio}

\subsection{El clítico \textit{se} como alomorfo de \textit{le}}\label{sec:10-alomorfo}

A partir de los resultados del ejercicio anterior, ¿cómo se puede explicar lo que
sucede en el ejemplo~5, en el que el pronombre que precede a \forma{lo}, con
función de COI, es \forma{se} en vez de \forma{le}? En este caso, \forma{se} no
es correferente del sujeto ---no remite a \forma{Elena}---, sino al COI \forma{a
su hermano} y, por lo tanto, no es un pronombre reflexivo, ya que se trata de
una variante de realización del clítico \forma{le}, es decir, de un
alomorfo\footnote{Este concepto se presentó en el \capref{4},
§\ref{sec:04-alomorfia}.}. Cuando \forma{le} precede a un clítico acusativo de
tercera persona (\forma{lo, la, los, las}), en su lugar aparece la forma
\forma{se}, que reemplaza a \forma{le} cuando tanto el COD como el COI son
pronominales; de modo que decimos \forma{Le di un regalo esta mañana}, pero
\forma{Se lo di esta mañana} y no \juicio{*}{\forma{Le lo di esta
mañana}}.\footnote{Según \textcite[173, n.~9]{ditullio2014}, «No se trata de un
condicionamiento fonológico, como lo prueba la existencia de secuencias como
\textit{le lavé la ropa}, \textit{le lotearon el terreno}». Se trata, pues, de un
condicionamiento morfológico propio de la concatenación entre los morfemas
pronominalizados de OD y OI.\pendiente{El original cita «Di Tullio (2010, 173,
nota 9)»}}

\section{Las construcciones pronominales}\label{sec:10-construcciones}

La aparición de los pronombres personales clíticos en la sintaxis no es un
fenómeno transparente y unívoco. Por el contrario, se trata de uno de los temas
más complejos y debatidos en las gramáticas. En primer lugar, puede haber más de
una forma para una misma función, como vimos en el ejercicio anterior, en el que
el dativo de tercera persona presenta dos variantes, \forma{se} y \forma{le},
cuya aparición depende del contexto oracional. En segundo lugar, una forma
pronominal no siempre se corresponde con una única función sintáctica, ya que,
como también vimos, puede manifestar caso acusativo o dativo, y cumplir función
de COD o de COI, según la oración. En tercer lugar, según veremos un poco más
adelante (§\ref{sec:10-cuasirreflejas}), los clíticos no siempre tienen función
sintáctica ni desempeñan papel temático.

\subsection{Prueba de sustitución}\label{sec:10-sustitucion}

Para reconocer la función sintáctica de los clíticos de primera y segunda
persona, que no poseen formas diferenciadas para acusativo y
dativo,\footnote{Para esta prueba se tienen en cuenta los datos propios de la
variedad del español rioplatense, en la que la pronominalización distingue entre
las formas \forma{lo, la, los, las} para el COD y las formas \forma{le, les}
para el COI. Existen variedades leístas, en las que se utilizan \forma{le, les}
tanto para el COI como para el COD (\forma{Le veo con frecuencia; Ese cuadro le
compramos el año pasado}). Asimismo, en las variedades laístas de ciertas zonas
del centro y norte de España, ocurre el fenómeno inverso: se utilizan las formas
\forma{la, las} para el COI femenino (\forma{A Luisa la dije que viniera; A la
pared la dieron una mano de pintura}) y, en las regiones loístas, se emplea
\forma{lo, los} para el COI masculino (\forma{A sus socios los mintió; Al vino
lo echan de todo}) \parencite{bosque2019}.} se puede aplicar una
\termino{prueba de sustitución}, mediante la cual se los reemplaza por un
pronombre de tercera persona, que tiene formas diferenciadas para acusativo y
dativo. Si al aplicar la prueba se cambia el clítico por \forma{la(s)} o
\forma{lo(s)}, se trata de un COD; en cambio, si se reemplaza por \forma{le(s)},
es un COI.

\begin{exe}
  \ex \textit{Me} encontró dormida. $\rightarrow$ \textit{La / *Le} encontró
      dormida. $\rightarrow$ \forma{Me} es COD.
  \ex \textit{Te} mostraron las fotos del viaje. $\rightarrow$ \textit{Le / *Lo}
      mostraron las fotos del viaje. $\rightarrow$ \forma{Te} es
      COI.\pendiente{El original dice «Me es COI»}
\end{exe}

\begin{ejercicio}\label{ej:10-sustitucion}
Según se muestra en el ejemplo resuelto~1:

\begin{consignas}
  \item aplique la prueba de sustitución para determinar qué función cumplen
        los pronombres en cursiva;
  \item luego, explique el resultado obtenido.
\end{consignas}

\begin{exe}
  \ex \textit{Me} arreglé frente al espejo.
  \ex \textit{Me} arreglé el pelo frente al espejo.
  \ex \textit{Se} vistió de gala.
  \ex ¿\textit{Te} ayudaron con la tarea?
  \ex \textit{Nos} dimos la vacuna para viajar.
\end{exe}

\begin{resuelto}
1. a) Sustitución por pronombre de tercera persona singular: \forma{Lo arreglé
frente al espejo}. b) El pronombre manifiesta caso acusativo; por lo tanto, su
función sintáctica es COD.
\end{resuelto}
\end{ejercicio}

\begin{ejercicio}\label{ej:10-papeles}
En los siguientes ejemplos analice el papel temático y la función sintáctica de
los pronombres en cursiva, según se muestra en el ejemplo resuelto~1.

\begin{exe}
  \ex \textit{Me} encontró dormida y \textit{te lo} contó.
  \ex \textit{Te} vimos cuando \textit{nos} comprabas los helados.
  \ex \textit{Se lo} llevé a su casa.
  \ex \textit{Me la} entregó en mi casa.
  \ex Los chicos \textit{se las} regalaron a sus amiguitos.
\end{exe}

\begin{resuelto}
1. \forma{Me}: paciente, COD; \forma{te}: meta, COI; \forma{lo}: tema, COD.
\end{resuelto}
\end{ejercicio}

\subsection{Transitividad y reflexividad}\label{sec:10-transitividad}

Como se planteó en el \capref{7}, la base de la \termino{transitividad} reside en
un sujeto que «actúa» (física o verbalmente) sobre un objeto, que siempre
representa una entidad diferente del sujeto.

\begin{exe}
  \ex {}[\etq{SN/Suj} El pastelero] preparó [\etq{SN/COD} varios postres] para la
      boda.
  \ex {}[\etq{SN/Suj} El automóvil] embistió [\etq{SN/COD} al peatón].
  \ex {}[\etq{SN/Suj} Los camilleros] trasladaron [\etq{SN/COD} al paciente].
\end{exe}

En \exref{ex:10-refl1} y \exref{ex:10-refl2} (que retoman \exref{ex:10-yomisma} y
\exref{ex:10-paciente2}), el panorama respecto de la transitividad cambia: sujeto
y objeto no solo coinciden en la persona gramatical, sino que corresponden al
mismo referente extralingüístico. Cuando la entidad a la que remite el objeto es
la misma que la del sujeto, hablamos de \termino{reflexividad}.

\begin{exe}
  \ex\label{ex:10-refl1} {}[\etq{1.ª pers.} Yo misma] [\etq{Pr Refl 1.ª pers.}
      \textit{me}] examiné las mamas cuidadosamente.
  \ex\label{ex:10-refl2} {}[\etq{3.ª pers.} La paciente] [\etq{Pr Refl 3.ª
      pers.} \textit{se}] examinó \textit{a sí misma} las mamas cuidadosamente.
\end{exe}

Además del pronombre reflexivo, las oraciones anteriores incluyen un COD que no
remite a alguien diferente, sino que es parte (\forma{las mamas}) del mismo
sujeto (\forma{yo} en \exref{ex:10-refl1}; \forma{la paciente} en
\exref{ex:10-refl2}), es decir, se trata de lo que se denomina una
\termino{posesión inalienable}.\footnote{\label{fn:10-inalienable}Posesión
inalienable es la que no puede ser separada de quien la posee porque es parte
integrante de él o de ella, como las partes del cuerpo o incluso, en algunos
casos, las pertenencias de su esfera personal.} En \exref{ex:10-refl1}, la
reflexividad se refuerza en el sujeto por el agregado del modificador enfático
\forma{misma} y, en \exref{ex:10-refl2}, mediante la forma \forma{a sí misma},
que duplica al COI.\pendiente{El original dice «a mí misma»}

\begin{recuadro}[Para ampliar: ¿qué sucede en las oraciones reflexivas con la transitividad?]
Aunque la forma transitiva se mantiene en la conformación estructural de estas
oraciones, que formalmente tienen un sujeto y un COD, ambos remiten a la misma
entidad del mundo extralingüístico, ya que el sujeto interviene en un evento que
revierte sobre sí mismo y no sobre un objeto externo a él. Por tanto, según
\textcite[499--500]{morenocabrera1991}, la reflexividad es «[u]no de los
procesos típicos que conduce a la intransitividad. La situación reflexiva típica
es aquella en la que el sujeto y el objeto son referencialmente idénticos
[\ldots] de modo que, aunque la oración sea sintácticamente transitiva, desde el
punto de vista semántico, se considera que es intransitiva».
\end{recuadro}

\section{Construcciones pronominales reflejas}\label{sec:10-reflejas}

Las \termino{construcciones pronominales} están formadas por un verbo y un
pronombre no oblicuo que coinciden en los rasgos flexivos de número y persona.
Dentro de estas construcciones, se consideran \termino{reflejas} aquellas en las
que, además de que la referencia del pronombre coincide con la del sujeto, el
pronombre representa uno de los argumentos del verbo y tiene la función
sintáctica de COD o de COI, según los casos, como ya hemos visto. Entre las
reflejas se distinguen dos grupos principales: \termino{reflexivas} y
\termino{recíprocas}. En las reflexivas, el sujeto interviene en un evento que
revierte sobre sí mismo (\forma{bañarse, peinarse}) o sobre alguna parte de sí
(\forma{cortarse las uñas, peinarse el flequillo}). En las recíprocas, el sujeto
es siempre plural y el evento implica una realización cruzada entre los
participantes (\forma{saludarse, darse las manos, abrazarse}). A continuación,
veremos las principales características de los distintos subgrupos.

\subsection{Construcciones pronominales reflexivas}\label{sec:10-reflexivas}

Dentro de las construcciones pronominales reflexivas distinguimos tres
subgrupos: reflexivas puras, reflexivas accidentales y reflexivas factitivas.

\subsubsection{Reflexivas puras}\label{sec:10-puras}

Son las reflejas más típicas y ocurren en todas las personas del paradigma
verbal. En ellas, el evento que realiza el sujeto recae sobre sí mismo
(\forma{bañarse, cortarse las uñas, rascarse}). El pronombre cumple la función
sintáctica de COD \exref{ex:10-marta} o de COI \exref{ex:10-sofia}. Como vimos en
\exref{ex:10-refl2}, como prueba de reconocimiento se puede incorporar el
sintagma con valor reflexivo \forma{a sí mismo/a}, que duplica el COD o el COI y
confirma que el pronombre es correferencial con el sujeto.

\begin{exe}
  \ex\label{ex:10-marta}
    \begin{xlist}
      \ex Marta \textit{se} baña con jabón biodegradable.
      \ex {}[\etq{SC=SF/O} [\etq{SN/Suj} Marta] [\etq{SV/Pred} [\etq{Pr Refl/COD}
          \textit{se}] baña\etq{V/n} [\etq{SN/COD} \textit{a sí misma}]
          [\etq{SP/ACI} con jabón biodegradable]]]
    \end{xlist}
  \ex\label{ex:10-sofia}
    \begin{xlist}
      \ex Sofía \textit{se} cortó las uñas esta mañana.
      \ex {}[\etq{SC=SF/O} [\etq{SN/Suj} Sofía] [\etq{SV/Pred} [\etq{Pr Refl/COI}
          \textit{se}] cortó\etq{V/n} [\etq{SN/COD} las uñas] [\etq{SN/COI}
          \textit{a sí misma}]]]
    \end{xlist}
\end{exe}

Típicamente, las reflexivas puras se construyen con verbos de acción que tienen
sujeto agente, \exref{ex:10-marta} y \exref{ex:10-sofia}, pero también admiten
verbos con sujeto experimentante, \exref{ex:10-adolescentes} y
\exref{ex:10-genio}:

\begin{exe}
  \ex\label{ex:10-adolescentes} Los adolescentes aún no se conocen bien a sí
      mismos.
  \ex\label{ex:10-genio} Ese escritor se cree un genio.
\end{exe}

\subsubsection{Reflexivas accidentales}\label{sec:10-accidentales}

En las reflejas que hemos estado considerando se presentan acciones realizadas
deliberadamente por los sujetos, que mayormente cumplen el papel temático de
agente. Pero veamos ahora qué sucede en \exref{ex:10-dedo}: ¿se trata también
de una acción deliberada?

\begin{exe}
  \ex\label{ex:10-dedo} Sofía \textit{se} cortó el dedo con un cuchillo.
\end{exe}

El ejemplo de \exref{ex:10-dedo} puede tener dos interpretaciones: la
«masoquista» o intencional, en la que Sofía se cortó un dedo voluntariamente, o
aquella en que el suceso le ocurrió en forma accidental. En el primer caso,
tenemos una construcción reflexiva pura, en la que Sofía es el agente que
deliberadamente realiza una acción que recae sobre ella misma. Sin embargo, la
interpretación «más natural» de esta oración es la segunda, en la que el corte
se ha producido de manera accidental, de modo que Sofía no interviene como
agente ejecutor, sino que, dado que el evento le sucede de modo involuntario, le
corresponde el papel temático de paciente.

Las construcciones que aceptan una interpretación de este tipo se conocen como
\termino{reflexivas accidentales} y presentan semejanzas y diferencias respecto
de las reflexivas puras. Al igual que las reflexivas puras
\exref{ex:10-paradigma-a}, las accidentales son paradigmáticas (es decir, de
toda persona) \exref{ex:10-paradigma-b} y, además de la construcción reflexiva
con un clítico de persona coincidente con el sujeto, también admiten clíticos
oblicuos (no coincidentes) con interpretación de acción deliberada o accidental
\exref{ex:10-oblicuo}.

\begin{exe}
  \ex\label{ex:10-paradigma}
    \begin{xlist}
      \ex\label{ex:10-paradigma-a} Reflexiva pura: \forma{Me peino. / Te peinás.
          / Se peina(n). / Nos peinamos\ldots}
      \ex\label{ex:10-paradigma-b} Reflexiva accidental: \forma{Me lastimé. / Te
          lastimaste. / Se lastimó (lastimaron). / Nos lastimamos\ldots}
    \end{xlist}
  \ex\label{ex:10-oblicuo}
    \begin{xlist}
      \ex Reflexiva pura: \forma{Me peino} (clítico coincidente con el sujeto) /
          \forma{Lo peino} (clítico no coincidente).
      \ex Reflexiva accidental: \forma{Me lastimé} (clítico coincidente) /
          \forma{Lo lastimé en la pelea} (clítico no coincidente con
          interpretación deliberada) / \forma{Lo lastimé sin darme cuenta}
          (clítico no coincidente con interpretación accidental).
    \end{xlist}
\end{exe}

Sin embargo, las accidentales también presentan diferencias respecto de las
reflexivas puras, algunas de las cuales las acercan a las cuasirreflejas medias,
que veremos más adelante (v.~§\ref{sec:10-medias}), como, por ejemplo, no tener
un sujeto agente ni admitir los refuerzos típicos de las reflexivas puras, como
\forma{a mí mismo/a, a vos mismo/a}, etc. Las reflexivas accidentales admiten
refuerzos del tipo de \forma{involuntariamente, accidentalmente} o \forma{por
casualidad}.

\begin{exe}
  \ex
    \begin{xlist}
      \ex {}[\etq{agente} Iván] se quebró (a sí mismo) una pierna para no ir a
          la guerra.\footnote{En este ejemplo, la presencia de un adjunto
          circunstancial de fin es un claro diagnóstico de la agentividad de
          Iván.}
      \ex\label{ex:10-ivan-b} {}[\etq{paciente} Iván] se quebró una pierna en el
          partido de rugby (*a sí mismo / accidentalmente).
    \end{xlist}
\end{exe}

Según vemos en \exref{ex:10-ivan-b}, aunque se acercan a las cuasirreflejas
medias, que dan cuenta de procesos no agentivos (\forma{El árbol se quebró por la
tormenta / solo}), las reflexivas accidentales no se intransitivizan, sino que
pueden presentar un COD referido a la parte del sujeto en que se produjo la
lesión, de modo que el papel temático que le corresponde es el de locativo, como
si dijéramos \sig{Sofía tuvo un corte en el dedo} o \sig{Juan sufrió una
fractura en la pierna}.

El carácter de reflexivas de las accidentales se manifiesta en que la acción,
aunque involuntaria, revierte sobre el sujeto y el clítico tiene función
sintáctica de COD \exref{ex:10-espina-a} o de COI \exref{ex:10-espina-b}.

\begin{exe}
  \ex
    \begin{xlist}
      \ex\label{ex:10-espina-a} Se pinchó con una espina de cactus.
      \ex\label{ex:10-espina-b} Se pinchó el dedo con una espina de cactus.
    \end{xlist}
\end{exe}

\subsubsection{Reflexivas factitivas}\label{sec:10-factitivas}

Comparemos estas dos oraciones:

\begin{exe}
  \ex\label{ex:10-selfi} Ayer \textit{me} saqué una selfi.
  \ex\label{ex:10-foto} El martes fui a un estudio fotográfico y \textit{me}
      saqué una foto para el registro de conducir.
\end{exe}

Aunque en ambas oraciones encontramos la construcción \forma{me saqué} referida a
una fotografía, solo la primera permite la inclusión del refuerzo \forma{a mí
mismo/a}; en la segunda, en cambio, es obvio que, si alguien va a un estudio
fotográfico, es para que otra persona le saque la foto. ¿Cómo podemos explicar
estas diferencias? En \exref{ex:10-selfi} encontramos las características de las
reflexivas puras que ya hemos visto en §\ref{sec:10-puras}, pero
\exref{ex:10-foto} implica que el sujeto no es el agente ejecutor de la acción,
que queda a cargo de otro ---la persona que en la casa de fotografía
efectivamente saca la foto---, quien es el verdadero agente, aunque este suele
quedar implícito o representado por un adjunto: \forma{Me saqué una foto con un
fotógrafo famoso}. Y ¿qué papel temático le corresponde al sujeto? Para
\textcite{kovacci1992}, quien promueve que otro realice la acción se considera un
\termino{causante instigador}. Las oraciones como \exref{ex:10-foto} se incluyen
con las reflexivas porque la acción que el causante instiga a otro a realizar
revierte sobre él mismo.

Las construcciones en las que el sujeto es distinto del ejecutor de la acción se
denominan \termino{reflexivas factitivas}. En ellas el pronombre cumple función
de COD o de COI. No admiten el refuerzo \forma{a mí mismo/a}, aunque aceptan una
paráfrasis con \forma{hacer}: \sig{Hice que otro (el fotógrafo) me sacara una
fotografía}.

\begin{exe}
  \ex {}[\etq{SC=SF/O} [\etq{Suj} \textit{pro}] [\etq{SV/Pred} [\etq{SN/ACT} el
      martes] [\etq{Pr Refl/COI} me] saqué\etq{V/n} [\etq{SN/COD} una foto]
      [\etq{SP/ACF} para el registro de conducir]]]
\end{exe}

Las factitivas son frecuentes en el uso cotidiano, tal como muestran los
siguientes ejemplos:

\begin{exe}
  \ex\label{ex:10-barberia} Hoy en la barbería \textit{me} corté el pelo. (=
      \sig{Hice que me cortaran el pelo})
  \ex Mi primo \textit{se} hizo una hermosa casa en un barrio privado de zona
      norte. (= \sig{Hizo que le hicieran / construyeran una hermosa casa\ldots})
  \ex \textit{Me} hice un piercing y ahora me sale pus de la nariz. (= \sig{Hice
      que me hicieran un piercing})
  \ex\label{ex:10-lorena} Lorena \textit{se} vacunó contra el sarampión esta
      mañana. (= \sig{Hizo que la vacunaran\ldots})
\end{exe}

\begin{ejercicio}\label{ej:10-factitivas}
En las oraciones anteriores, \exref{ex:10-barberia} a \exref{ex:10-lorena},
identifique el agente y el causante instigador. Señale la función sintáctica que
cumple el pronombre, según se muestra en el ejemplo resuelto de abajo.

\begin{resuelto}
\exref{ex:10-barberia}: agente: el barbero; causante instigador: el sujeto
tácito \forma{yo}. El pronombre \forma{me} es COI.\pendiente{El original
intercambia agente y causante instigador, también en el solucionario}
\end{resuelto}
\end{ejercicio}

\subsection{Construcciones pronominales recíprocas}\label{sec:10-reciprocas}

Las construcciones recíprocas responden a la caracterización de las reflejas: el
pronombre es correferencial con el sujeto y cumple función sintáctica de COD
\exref{ex:10-matias} o de COI \exref{ex:10-novios}.

\begin{exe}
  \ex\label{ex:10-matias} Matías y Francisco \textit{se} esperaron a la salida
      del colegio.
  \ex\label{ex:10-novios} Durante la ceremonia, los novios \textit{se} pusieron
      los anillos.
\end{exe}

Ahora bien, ¿qué diferencia estas construcciones de las reflexivas? Aquí también
el evento «transita» de un sujeto a un objeto que remite al mismo referente, pero
se diferencian porque se trata de acciones cruzadas: en \exref{ex:10-matias},
Matías esperó a Francisco y, recíprocamente, Francisco esperó a Matías, y, en
\exref{ex:10-novios}, la novia le puso el anillo al novio mientras él, a su vez,
se lo colocaba a ella. Los eventos recíprocos ocurren simultáneamente (o casi),
ya que la acción realizada por cada agente revierte en otro, que hace las veces
de paciente \exref{ex:10-recip-a} o meta \exref{ex:10-recip-b}.

\begin{exe}
  \ex
    \begin{xlist}
      \ex\label{ex:10-recip-a} {}[\etq{SC=SF/O} [\etq{SN/Suj} Matías y Francisco]
          [\etq{SV/Pred} [\etq{Pr Recipr/COD} se] esperaron\etq{V/n}
          [\etq{SP/ACT} a la salida del colegio]]]
      \ex\label{ex:10-recip-b} {}[\etq{SC=SF/O} [\etq{SN/Suj} los novios]
          [\etq{SV/Pred} [\etq{Pr Recipr/COI} se] pusieron\etq{V/n} [\etq{SN/COD}
          los anillos]]]
    \end{xlist}
\end{exe}

Como se trata de una acción que va de uno a otro, estas construcciones siempre
exigen un sujeto que exprese pluralidad, ya sea morfológicamente
\exref{ex:10-novios}, sintácticamente, mediante la coordinación de SSNN
\exref{ex:10-matias}, o semánticamente, por medio de un sustantivo colectivo
(\forma{La tripulación se abrazaba feliz por haber llegado a destino}).

Para reconocer estas construcciones se pueden aplicar las siguientes pruebas:

\begin{itemize}
  \item agregar adverbios como \forma{mutuamente} o \forma{recíprocamente}, el SP
        \forma{entre sí} o construcciones como \forma{uno al otro}:
\end{itemize}

\begin{exe}
  \ex {}[\etq{SC=SF/O} [\etq{SN/Suj} los novios] [\etq{SV/Pred} [\etq{Pr
      Recipr/COI} se] pusieron\etq{V/n} [\etq{SN/COD} los anillos]
      [\etq{SAdv/ACM} mutuamente]]]\pendiente{El original rotula aquí el
      pronombre como COD}
\end{exe}

\begin{itemize}
  \item dado que el sujeto es plural, también se pueden descomponer sus miembros
        en dos oraciones coordinadas por \forma{y}:
\end{itemize}

\begin{exe}
  \ex Matías esperó a Francisco y Francisco esperó a Matías.
\end{exe}

\begin{itemize}
  \item pronominalizar separadamente los dos OODD:
\end{itemize}

\begin{exe}
  \ex Matías \textit{lo} esperó a Francisco y Francisco \textit{lo} esperó a
      Matías.\pendiente{El original dice «Federico» en estas dos pruebas}
\end{exe}

\begin{destacado}
\textbf{Para tener en cuenta.} Cuando el sujeto es plural, se puede producir
ambigüedad entre la interpretación reflexiva y la recíproca, y será el contexto
el que permita desambiguar. Así, en \forma{Luisa y Mara se miraron en el espejo
antes de entrar en el salón}, podemos interpretar que \sig{cada una se miró a sí
misma en el espejo} (reflexiva pura) o que \sig{a través del espejo, una miró a
la otra, es decir, que se miraron mutuamente} (recíproca). El contexto suele
contener «pistas» que nos permiten desambiguar ambas lecturas:

\begin{exe}
  \ex
    \begin{xlist}
      \ex Luisa y Mara se miraron en el espejo y cada una revisó su peinado.
          (reflexiva pura)
      \ex Luisa y Mara se miraron en el espejo y ambas se sonrieron en forma
          cómplice. (recíproca)
    \end{xlist}
\end{exe}
\end{destacado}

\begin{ejercicio}\label{ej:10-ambiguas}
Teniendo en cuenta que las siguientes oraciones son ambiguas (admiten más de una
interpretación), resuelva lo pedido más abajo.

\begin{exe}
  \ex Ambos luchadores \textit{se torcieron} la muñeca durante la competencia.
  \ex Ana y su hermano \textit{se tiñeron} el pelo.
  \ex Los dos estudiantes \textit{se copiaron} en el examen.
\end{exe}

\begin{consignas}
  \item Provea dos contextos que permitan desambiguarlas, según se muestra en el
        ejemplo resuelto~1.
  \item En 2, además de la reflexiva pura y la recíproca, ¿qué otra
        interpretación es posible?
\end{consignas}

\begin{resuelto}
1. Reflexiva pura: \forma{Uno se torció su propia muñeca al caer y el otro la
suya, contra las sogas}. Recíproca: \forma{Se torcieron uno al otro la muñeca
como resultado de una toma peligrosa en la que ambos participaron}.
\end{resuelto}
\end{ejercicio}

\begin{ejercicio}\label{ej:10-tiporefleja}
Analice los siguientes ejemplos y decida de qué tipo de refleja se trata en cada
caso.

\begin{exe}
  \ex Juanita \textit{se} vistió de rojo para la fiesta y \textit{se} pintó los
      ojos de dorado.
  \ex Silvina decidió tatuar\textit{se} el nombre de su hijo en el brazo.
  \ex \textit{Se} depiló las cejas para emprolijar su mirada.
  \ex El pintor holandés Vincent van Gogh \textit{se} cortó una oreja.
  \ex Hipólito Yrigoyen y Lisandro de la Torre \textit{se} enfrentaron en un
      duelo.
\end{exe}

\begin{resuelto}
1. Ambos pronombres \forma{se} forman construcciones reflexivas puras.
\end{resuelto}
\end{ejercicio}

\begin{ejercicio}\label{ej:10-danna}
Analice todas las construcciones subrayadas en el siguiente fragmento;
identifique las que son pronominales y aporte pruebas que sustenten sus
respuestas, según se muestra en el ejemplo resuelto de abajo.

% Texto periodístico protegido: se incorpora a mano desde el original.
\textoaincorporar{el fragmento adaptado de una nota de \textit{Infobae}
(14/04/2020) sobre Danna Paola, con las construcciones subrayadas: \textit{se
había bañado; se toca el cabello; le arregle las pestañas y las cejas; me estoy
agarrando el pelo; me vean bien el pelo; me las tengo que hacer yo; me las haga;
se comió todas las uñas de la mano derecha}}

\begin{resuelto}
\forma{Se había bañado}: construcción pronominal refleja. El pronombre \forma{se}
es correferencial con el sujeto. Admite el refuerzo \forma{a sí misma}.
\end{resuelto}
\end{ejercicio}

\section{Pronominales sí, pero no reflexivas: las construcciones pronominales cuasirreflejas}\label{sec:10-cuasirreflejas}

No todas las construcciones pronominales son reflexivas: hay un grupo en las que
el verbo y el pronombre coinciden en número y persona, pero no comparten el resto
de las características de las reflexivas. Para ir diferenciando las reflexivas
de las otras, a las que \textcite{bello2002} denominó \termino{cuasirreflejas},
comencemos observando las siguientes frases extraídas a) de una pintada callejera
y b) de un afiche del gobierno de la provincia de Buenos Aires.

\imagenpendiente{a) pintada callejera; b) afiche del gobierno de la provincia de
Buenos Aires}

\begin{exe}
  \ex\label{ex:10-feliz} No sabía qué ponerme y \textit{me} puse feliz.
  \ex\label{ex:10-turnos} \textit{Se} están enviando nuevos turnos de vacunación.
  \ex\label{ex:10-vacunate} Buenos Aires, vacuná\textit{te}.
\end{exe}

Aunque todas las construcciones pronominales están formadas por un verbo y un
pronombre que comparten los rasgos flexivos de número y persona, las reflejas,
que hemos visto en §\ref{sec:10-reflejas}, deben cumplir otros dos requisitos
más: el pronombre debe desempeñar algún papel temático y la función sintáctica
de COD o de COI. En cambio, en las \termino{cuasirreflejas}, como las que
aparecen en las imágenes de arriba, aunque verbo y pronombre comparten los rasgos
de flexión, no cumplen los otros dos requisitos: el pronombre no expresa ningún
papel temático y tampoco tiene función sintáctica. Para
\textcite[§\,759]{bello2002}, a diferencia de las reflexivas, en las que agente y
paciente coinciden, en las cuasirreflejas «gramaticalmente parece decirse que el
sujeto obra en sí mismo», lo que significa que las construcciones cuasirreflejas
refieren a un proceso en el que el sujeto no es el agente que lo realiza; por lo
tanto, los refuerzos que habíamos visto con las reflexivas,
\exref{ex:10-desvesti} y \exref{ex:10-odian}, no se aplican,
\exref{ex:10-despidio} y \exref{ex:10-jactaban}:

\begin{exe}
  \ex\label{ex:10-desvesti} Cuando llegué a mi casa, \textit{me} desvestí (a mí
      misma) y \textit{me} di una ducha (a mí misma).
  \ex\label{ex:10-odian} Mi prima y su hermana \textit{se} odian (mutuamente /
      una a la otra).
  \ex\label{ex:10-despidio} Ya era tarde cuando Karina \textit{se} despidió (*a
      sí misma) de nosotros.
  \ex\label{ex:10-jactaban} Ambos \textit{se} jactaban de sus hazañas (*a sí
      mismos / *mutuamente).
\end{exe}

Las construcciones cuasirreflejas son sumamente complejas porque hay distintos
subgrupos cuyos límites «son inestables y, por decirlo así, porosos»
\parencite[17]{sanchezlopez2002}; por tanto, es difícil clasificarlas siguiendo
criterios que las abarquen a todas. En líneas generales, se trata de
construcciones que producen algún «cambio» en la oración. Pero ¿de qué tipo de
cambio se trata y a qué afecta? En las cuasirreflejas, el clítico no cumple
función sintáctica ni de COD ni de COI y se interpreta como una marca de función
que recibe el nombre de \termino{índice pronominal cuasirreflejo} (IPCR). La
introducción de este pronombre incide en la conformación de alguna de las tres
estructuras del predicado que presentamos en el \capref{6}: la estructura
argumental (EA), la estructura sintáctica (ES) o la estructura eventiva (EE), por
lo que dividiremos estas construcciones en tres grupos:

\begin{itemize}
  \item Construcciones cuasirreflejas que producen cambio en la estructura
        argumental. La introducción del clítico elimina la manifestación de
        algún argumento de la estructura sintáctica. Dentro de estas
        construcciones se encuentran las cuasirreflejas medias, las
        medio-pasivas, las pasivas y las impersonales.
  \item Construcciones cuasirreflejas que producen cambio en la estructura
        sintáctica. El clítico produce una modificación en la forma en que se
        manifiesta la transitividad. Las construcciones de este tipo son las
        cuasirreflejas alternantes y las inherentes.
  \item Construcciones cuasirreflejas que producen cambio en la estructura
        eventiva. La presencia del clítico ocasiona cambios en la aspectualidad.
        Se ubican en este grupo las construcciones cuasirreflejas que delimitan
        o intensifican el evento.
\end{itemize}

\subsection{Construcciones cuasirreflejas que producen cambio en la estructura argumental}\label{sec:10-ea}

Las cuasirreflejas que cambian la estructura argumental comprenden cuatro
subgrupos: medias, medio-pasivas, pasivas e impersonales.

Todas estas construcciones tienen en común que alteran la manifestación
sintáctica del argumento agente. Esto, en las construcciones medias,
medio-pasivas y pasivas, produce la promoción del tema/paciente a la función de
sujeto e intransitiviza la oración. En las impersonales, en cambio, ningún
argumento se promueve como sujeto y, como veremos, no todas las construcciones
se intransitivizan. Otra diferenciación importante entre estos subgrupos es que
solo las construcciones medias son paradigmáticas, es decir, el pronombre puede
aparecer en cualquiera de las personas gramaticales (primera, segunda o tercera)
y tanto en singular como en plural (\forma{Me entristecí / Te entristeciste / Se
entristecieron con las noticias}, etc.). En las medio-pasivas, pasivas e
impersonales, \forma{se} no es paradigmático: está fijado en tercera persona,
singular o plural en las medio-pasivas (\forma{Este error se subsana fácilmente}
/ \forma{Estos errores se subsanan\ldots}) y en las pasivas (\forma{Se presentó
un informe} / \forma{Se presentaron varios informes}), y en tercera persona
singular en las impersonales (\forma{Aquí no se fuma} /
\juicio{*}{\forma{Aquí no se fuman}}). En lo que sigue describiremos las
características de cada subgrupo.

\subsubsection{Construcciones cuasirreflejas medias}\label{sec:10-medias}

En las construcciones de este subgrupo se elimina el agente, lo que motiva la
promoción o ascenso del argumento tema/paciente, que se convierte en sujeto de
la oración.

\begin{exe}
  \ex {}[\etq{agente} Pedro] abrió [\etq{tema} la puerta]. $\rightarrow$
      [\etq{tema} La puerta] se abrió (sola).
\end{exe}

Como en todas las cuasirreflejas, en estas construcciones el pronombre no
desempeña ninguna función sintáctica y se considera una marca de función.
Asimismo, se destaca que el sujeto «sufre un proceso que escapa a su control
efectivo» \parencite[77]{sanchezlopez2002}; de ahí que se considere que estas
oraciones se corresponden con la \termino{voz media}.

\begin{destacado}
Según \textcite[1636]{mendikoetxea1999}, la voz media «se interpreta como la
forma de expresar que la acción que denota el verbo “afecta”, en mayor o menor
grado, al sujeto (incide sobre el sujeto, interesa al sujeto, indica un cambio en
el sujeto, etc.). En otras palabras, las construcciones medias se caracterizan
por tener sujetos “afectados”».
\end{destacado}

Este tipo de construcciones ocurre con verbos transitivos de \sig{cambio de
estado}, que se intransitivizan por la pronominalización. Según el tipo de cambio
involucrado, se reconocen tres subgrupos: con verbos que indican \sig{cambio de
estado físico} (\forma{romperse, secarse, hundirse}), con verbos que indican
\sig{cambio de estado mental} (\forma{entristecerse, asombrarse, divertirse}) y
con verbos que indican \sig{ubicación o posición} (\forma{acercarse, moverse,
levantarse}).

\paragraph{Medias con verbos de cambio de estado físico.} Prestemos atención a
los siguientes pares de oraciones:

\begin{exe}
  \ex\label{ex:10-pared}
    \begin{xlist}
      \ex\label{ex:10-pared-a} El albañil rompió la pared.
      \ex\label{ex:10-pared-b} La pared se rompió.
    \end{xlist}
  \ex\label{ex:10-ropa}
    \begin{xlist}
      \ex\label{ex:10-ropa-a} El viento secó la ropa.
      \ex\label{ex:10-ropa-b} La ropa se secó con el viento.
    \end{xlist}
\end{exe}

Tanto \exref{ex:10-pared-a} como \exref{ex:10-ropa-a} son oraciones transitivas
en las que el primer argumento es un agente (\forma{el albañil}) o una causa
(\forma{el viento}), y el segundo es un tema (\forma{la pared} o \forma{la
ropa}). Pero ¿qué sucede en \exref{ex:10-pared-b} y \exref{ex:10-ropa-b}? El tema
ha sido promovido a la función más prominente en la estructura oracional, la de
sujeto, lo que en \exref{ex:10-pared-b} produce la eliminación total del agente
(\juicio{*}{\forma{La pared se rompió por el albañil}}) o, como sucede en
\exref{ex:10-ropa-b}, la causa es relegada a una posición marginal, en la que
cumple la función de adjunto circunstancial.

¿Qué efecto producen estos cambios? En \exref{ex:10-belgrano} se muestra el
contraste entre una oración activa, de lectura agentiva \exref{ex:10-belgrano-a},
y una oración media \exref{ex:10-belgrano-b}, en la que la eliminación del agente
(o su democión a una función marginal) promueve una \termino{lectura de proceso}
que, en algunos casos, permite la introducción del refuerzo \forma{solo/a} o
\forma{por sí solo/a}.

\begin{exe}
  \ex\label{ex:10-belgrano}
    \begin{xlist}
      \ex\label{ex:10-belgrano-a} Los ingleses hundieron el \textit{General
          Belgrano} durante la guerra de Malvinas.\\
          {}[\etq{SC=SF/O} [\etq{SN/Suj/agente} los ingleses] [\etq{SV/Pred}
          hundieron\etq{V/n} [\etq{SN/COD/tema} el \textit{General Belgrano}]
          [\etq{SP/ACT} durante la guerra de Malvinas]]]
      \ex\label{ex:10-belgrano-b} El \textit{Titanic} se hundió (solo) / (por un
          iceberg / porque chocó con un iceberg).\\
          {}[\etq{SC=SF/O} [\etq{SN/Suj/tema} el \textit{Titanic}] [\etq{SV/Pred}
          [\etq{Pr/IPCR} se] hundió\etq{V/n} [\etq{SA/ACM} solo] / [\etq{SP/ACC}
          por un iceberg]]]
    \end{xlist}
\end{exe}

Los verbos que intervienen en estas construcciones se pueden caracterizar como
\sig{verbos de cambio} o que producen un cambio en el objeto. Muchos permiten lo
que se conoce como \termino{alternancia causativa} y ofrecen dos posibles
manifestaciones:

\begin{enumerate}[label=\alph*)]
  \item la \termino{construcción transitiva causativa} típica, en la que el
        cambio en el objeto es producido por un agente, \exref{ex:10-pared-a} y
        \exref{ex:10-belgrano-a}, o por una causa \exref{ex:10-ropa-a};
  \item la \termino{construcción media anticausativa}, en la que el agente se
        elimina \exref{ex:10-pared-b} o se demueve a una función sintáctica
        opcional, \exref{ex:10-ropa-b} y \exref{ex:10-belgrano-b}, y el tema se
        promueve a la función de sujeto. Esta transformación, que elimina al COD,
        produce una intransitivización o inacusativización, por lo que se dice
        que el pronombre cumple función intransitivizadora.
\end{enumerate}

\begin{recuadro}[Para ampliar]
No todos los verbos de cambio transitivos entran en la alternancia causativa.
Para que la alternancia sea posible, el evento tiene que poder ser interpretado
como un proceso que le ocurre al tema sin intervención de un agente. Así, verbos
transitivos típicamente agentivos como \forma{cortar, construir, leer,
asesinar}, por lo general, no admiten la alternancia, ya que siempre requieren un
sujeto agente.

\begin{exe}
  \ex
    \begin{xlist}
      \ex El pastelero cortó el postre en porciones.
      \ex[*]{El postre se cortó (solo) en porciones.}
    \end{xlist}
  \ex
    \begin{xlist}
      \ex Niemeyer construyó los principales edificios de Brasilia.
      \ex[*]{Los principales edificios de Brasilia se construyeron (solos).}
    \end{xlist}
\end{exe}

Por otra parte, los verbos de proceso tampoco se comportan todos igual. Para
explicar estas diferencias, \textcite{mendikoetxea2011} compara \forma{romper} y
\forma{estallar}, dos verbos que implican \sig{la desintegración material de un
objeto}, que representa al tema afectado por el evento. La autora encuentra que,
a pesar de su aparente similitud semántica, las construcciones que ambos verbos
permiten no son las mismas:

\begin{exe}
  \ex
    \begin{xlist}
      \ex La tormenta rompió el vidrio.
      \ex[*]{La tormenta estalló el vidrio.}
    \end{xlist}
  \ex
    \begin{xlist}
      \ex[*]{El vidrio rompió.}
      \ex El vidrio estalló.
    \end{xlist}
  \ex
    \begin{xlist}
      \ex El vidrio se rompió.
      \ex[*]{El vidrio se estalló.}
    \end{xlist}
\end{exe}

¿Cómo se pueden explicar estas diferencias? Si bien tanto \forma{romper} como
\forma{estallar} son verbos de cambio de estado, se diferencian porque pertenecen
a dos subclases distintas, según el evento sea producido por una causa externa o
por una causa interna:

\begin{enumerate}[label=\alph*)]
  \item \termino{Verbos de causa externa}, como \forma{romper, secar, quemar} o
        \forma{abrir}: refieren una acción que tiene una causa que la produce.
        Estos verbos admiten que el tema afectado por el evento aparezca tanto en
        la posición de objeto como en la de sujeto, pero en este segundo caso
        requieren la introducción del clítico \forma{se}:
        \begin{exe}
          \ex
            \begin{xlist}
              \ex Juana secó la ropa.
              \ex La ropa \textit{se} secó.
            \end{xlist}
          \ex
            \begin{xlist}
              \ex El incendio está quemando el bosque.
              \ex El bosque \textit{se} está quemando (debido al incendio).
            \end{xlist}
        \end{exe}
  \item \termino{Verbos de causa interna}, como \forma{estallar, arder} o
        \forma{envejecer}: no involucran causa ni agente externo, es decir que
        nadie o nada ocasiona el evento, sino que este simplemente se da u ocurre
        por un proceso interno, que se produce o tiene lugar en el objeto. Estos
        verbos solo admiten la construcción con el tema en función de sujeto y
        rechazan \forma{se}:
        \begin{exe}
          \ex
            \begin{xlist}
              \ex La leña verde no arde bien.
              \ex[*]{La leña verde no \textit{se} arde bien.}
            \end{xlist}
          \ex
            \begin{xlist}
              \ex Algunas personas envejecen muy bien.
              \ex[*]{Algunas personas \textit{se} envejecen muy bien.}
            \end{xlist}
        \end{exe}
\end{enumerate}
\end{recuadro}

\paragraph{Medias con verbos de cambio de estado mental (verbos psicológicos).}
Los verbos psicológicos, como \forma{entristecerse, enojarse, divertirse,
sorprenderse}, también se incluyen entre los de cambio de estado y señalan un
cambio interno de tipo mental que alguien experimenta ---\forma{la abuela}
\exref{ex:10-abuela-a} y \forma{la población} \exref{ex:10-abuela-b}---.

\begin{exe}
  \ex
    \begin{xlist}
      \ex\label{ex:10-abuela-a} El nieto alegró a su abuela con una sonrisa.\\
          {}[\etq{SC=SF/O} [\etq{SN/Suj/agente} el nieto] [\etq{SV/Pred}
          alegró\etq{V/n} [\etq{SN/COD/experimentante} a su abuela] [\etq{SP/ACI}
          con una sonrisa]]]
      \ex\label{ex:10-abuela-b} Un nuevo virus preocupa a la población.
    \end{xlist}
\end{exe}

Al igual que los verbos de cambio de estado físico, los que indican cambio
psicológico admiten alternancia entre:

\begin{enumerate}[label=\alph*)]
  \item la construcción transitiva, con agente \exref{ex:10-abuela-a} o causa
        \exref{ex:10-abuela-b} y un experimentante;
  \item la construcción media, en la que el agente o la causa se elimina o
        demueve a una función sintáctica opcional de adjunto y el experimentante
        es promovido a la función de sujeto mediante un proceso de
        intransitivización.
\end{enumerate}

\begin{exe}
  \ex
    \begin{xlist}
      \ex La abuela se alegró con la sonrisa de su nieto.\\
          {}[\etq{SC=SF/O} [\etq{SN/Suj/experimentante} la abuela]
          [\etq{SV/Pred} [\etq{Pr/IPCR} se] alegró\etq{V/n} [\etq{SP/ACI} con la
          sonrisa de su nieto]]]
      \ex La población se preocupó por el nuevo virus.
    \end{xlist}
\end{exe}

\paragraph{Medias con verbos de cambio de ubicación o posición.} Dentro del
amplio grupo de verbos de movimiento, en este apartado nos centramos en los que
son transitivos y admiten la alternancia causativa. Más adelante, dentro de las
construcciones cuasirreflejas que producen cambio en la estructura eventiva,
tratamos algunos verbos de movimiento intransitivos que presentan una variante
pronominal (v.~§\ref{sec:10-delimitacion}).

Si bien la mayoría de los autores incluyen estos verbos con los de cambio
físico, los de ubicación o posición presentan características que los
distinguen. Cotejemos los siguientes ejemplos para ver sus particularidades.

\begin{exe}
  \ex
    \begin{xlist}
      \ex\label{ex:10-lobo-a} Los cazadores alejaron al lobo.
      \ex El lobo se alejó.
    \end{xlist}
  \ex\label{ex:10-sento}
    \begin{xlist}
      \ex\label{ex:10-sento-a} La enfermera sentó al paciente.
      \ex\label{ex:10-sento-b} El paciente se sentó por sí solo.
    \end{xlist}
  \ex
    \begin{xlist}
      \ex\label{ex:10-chofer-a} El chofer detuvo el coche.
      \ex El coche se detuvo (porque el chofer paró el motor / porque se paró el
          motor) $\rightarrow$ el motor del automóvil se detuvo.
    \end{xlist}
\end{exe}

En casos como los anteriores, con verbos transitivos, las construcciones
cuasirreflejas con verbos de cambio de ubicación o posición se comportan como las
de cambio físico o mental que ya hemos visto, es decir que alternan entre:

\begin{enumerate}[label=\alph*)]
  \item una construcción transitiva causativa, con sujeto agente (\forma{los
        cazadores, la enfermera, el chofer}) y un objeto tema o paciente
        (\forma{el lobo, el paciente, el coche});
  \item una construcción media anticausativa, formada mediante la inclusión de
        \forma{se}, en la que se elimina al agente, promoviendo al tema o
        paciente a la función de sujeto.
\end{enumerate}

Pero ¿qué particularidad presenta el nuevo sujeto de la construcción pronominal?
¿Es un tema/paciente como en los verbos de cambio de estado? Observemos el
contraste que presenta \exref{ex:10-muneca} respecto de \exref{ex:10-sento}.

\begin{exe}
  \ex\label{ex:10-muneca}
    \begin{xlist}
      \ex\label{ex:10-muneca-a} La niña sentó a su muñeca.
      \ex[?]{\label{ex:10-muneca-b} La muñeca se sentó (sola).}
    \end{xlist}
\end{exe}

La construcción transitiva de \exref{ex:10-muneca-a} es paralela a la de
\exref{ex:10-sento-a}; sin embargo, aunque podríamos pensar en una situación
particular en que \exref{ex:10-muneca-b} fuera aceptable, fuera de contexto
resulta anómala. Para \textcite[200]{kovacci1992}, en estas construcciones,
cuando el sujeto es animado, como en \exref{ex:10-sento-b}, realiza
simultáneamente al agente y al paciente, ya que «[\ldots] dos participantes
pueden ser correferenciales y manifestarse en la misma función». También para
\textcite[1593]{mendikoetxea1999}, dado que estos son verbos de causa externa,
«el sujeto se interpreta como causa y tema»,\footnote{Para
\textcite[495]{martinezrojas2021}, a diferencia de las anticausativas, que
refieren eventos espontáneos con sujeto paciente, las que se forman con verbos
de movimiento del cuerpo (\textit{body motion verbs}) (autocausativas o
endorreflexivas) refieren a «eventos en los que el sujeto animado es agente».}
ya que podemos tener contrastes como \forma{Yo me levanto solo (voz media) y no
espero que otros me levanten (voz activa)} frente a \forma{Yo me levanto solo y
no espero ser levantado (por otros) (voz pasiva)}. Asimismo, estas oraciones
admiten la construcción final con \forma{para}, que se considera un diagnóstico
de agentividad: \forma{Yo me levanto para no seguir escuchando lo que dicen / El
paciente se sentó para comer / El cazador se alejó para perseguir al lobo}.

Y ¿qué encontramos en oraciones como \exref{ex:10-tormenta} a
\exref{ex:10-findeano}, en las que estos verbos se construyen con sujetos no
animados? Se trata de entidades que por sí solas desencadenan un proceso. En
\exref{ex:10-findeano}, aunque el sujeto tácito \forma{nosotros} sea
\sig{animado} y \sig{humano}, no posee voluntariedad para realizar un evento como
\sig{acercarse a fin de año}, que se presenta como un proceso vinculado con el
transcurso temporal. En estos casos, los sujetos solo manifiestan el papel
temático de tema o paciente, según corresponda.

\begin{exe}
  \ex\label{ex:10-tormenta} La tormenta se aproxima a la ciudad.
  \ex La solución se asentó en el fondo del recipiente.
  \ex\label{ex:10-findeano} Nos acercamos a fin de año.
\end{exe}

\begin{destacado}
\textbf{Para tener en cuenta.} Las diferencias presentadas a través de los
ejemplos incluidos en este apartado muestran que los verbos tienen siempre
múltiples posibilidades de construcción y de combinación con diferentes
complementos (incluido el sujeto) y adjuntos y que, por eso, la interpretación de
cada oración debe hacerse teniendo en cuenta todos estos factores.
\end{destacado}

\begin{ejercicio}\label{ej:10-levantar}
Compare las oraciones de los siguientes grupos y resuelva lo pedido, según se
muestra en el ejemplo resuelto~1.

\begin{consignas}
  \item Señale semejanzas y diferencias.
  \item Aplique alguna prueba para corroborar la agentividad de los sujetos.
  \item ¿Cuál es el papel temático de \forma{las sillas} en 1a y en 1c?
  \item ¿Qué refuerzo puede aplicar en 2b?
\end{consignas}

\begin{exe}
  \ex
    \begin{xlist}
      \ex Los asistentes levantaron las sillas en señal de protesta.
      \ex Los asistentes levantaron la mano para hacer preguntas.
      \ex Los asistentes se levantaron de sus sillas para aplaudir de pie al
          cantante.
    \end{xlist}
  \ex
    \begin{xlist}
      \ex El nuevo gobierno extendió los límites de su dominio.
      \ex El incendio se extendió por el bosque nativo.
    \end{xlist}
\end{exe}

\begin{resuelto}
1. Semejanzas entre 1a y 1b: en ambas, \forma{levantar} tiene empleo transitivo
con sujeto agente ---\forma{los asistentes}--- y COD tema ---\forma{las sillas}
(1a) y \forma{la mano} (1b)---, y aceptan la paráfrasis con \forma{hacer}:
\forma{¿Qué hicieron los asistentes?} Diferencia entre 1a y 1b: en 1b el COD es
una posesión inalienable, \forma{la mano}. La diferencia con 1c es que se trata
de una construcción media, en la que el sujeto es a la vez agente y paciente. El
papel temático de \forma{las sillas} en 1a es tema y en 1c es
origen.\pendiente{El original dice «en (b)» en la última oración}
\end{resuelto}
\end{ejercicio}

\begin{ejercicio}\label{ej:10-cambio}
Compare las estructuras de cada par, en las que el verbo aparece con y sin
pronombre. Luego resuelva lo pedido en las consignas, según se muestra en el
ejemplo resuelto~1.

\begin{consignas}
  \item Identifique la clase del verbo según el tipo de cambio que produce (de
        estado físico, de estado mental, de ubicación o posición).
  \item Teniendo en cuenta la estructura argumental de los verbos, trate de
        explicar qué cambios produce la introducción del pronombre.
\end{consignas}

\begin{exe}
  \ex
    \begin{xlist}
      \ex Juan abrió la puerta con la llave.
      \ex Se abrió la puerta con el viento.
    \end{xlist}
  \ex
    \begin{xlist}
      \ex Levanté las cortinas para que entre la luz.
      \ex Me levanto temprano y estudio.
    \end{xlist}
  \ex
    \begin{xlist}
      \ex Acuesto a los niños luego de bañarlos.
      \ex Se acostó mientras los niños dormían.
    \end{xlist}
\end{exe}

\begin{resuelto}
1. Verbo de cambio de estado físico. 1a: construcción transitiva con agente y
tema (\forma{la puerta}). 1b: construcción cuasirrefleja con promoción del tema
(\forma{la puerta}) a sujeto y democión de la causa a ACC (\forma{con el
viento}).
\end{resuelto}
\end{ejercicio}

\subsubsection{Construcciones cuasirreflejas medio-pasivas}\label{sec:10-mediopasivas}

Al igual que en las construcciones medias, también en las medio-pasivas se
demueve o elimina el agente, pero, a diferencia de lo que sucede con el grupo
anterior, en las medio-pasivas el agente eliminado tiene interpretación
indefinida. Se trata de oraciones como las siguientes:

\begin{exe}
  \ex\label{ex:10-problema} El problema de matemáticas se soluciona en dos
      minutos / en dos patadas.
  \ex\label{ex:10-libros} Los libros de divulgación científica se leen
      fácilmente / rápidamente / de corrido.
  \ex\label{ex:10-pastillas} Estas pastillas se toleran bien / sin problemas.
\end{exe}

Una lectura atenta de estas oraciones pone en evidencia que comparten
características con las pasivas (v.~§\ref{sec:10-pasivas}) y con las
impersonales (v.~§\ref{sec:10-impersonales}) con \forma{se}: en primer lugar,
solo se pueden construir en tercera persona, es decir que el único pronombre que
aparece es \forma{se}. En segundo lugar, como las pasivas, admiten concordancia
en singular \exref{ex:10-problema} y en plural, \exref{ex:10-libros} y
\exref{ex:10-pastillas}, y, en ambas construcciones, el tema o paciente se
promueve a la posición de sujeto. También como las pasivas, estas construcciones
implican la eliminación del agente de la correspondiente oración activa, pero se
diferencian porque en las medio-pasivas el agente eliminado tiene interpretación
indefinida: \forma{Este problema se soluciona en dos minutos} $\rightarrow$
\sig{Cualquiera lo soluciona / lo puede solucionar}. Otra característica propia
de estas construcciones es que, a diferencia de las pasivas y de las
impersonales, que son eventivas, las medio-pasivas constituyen predicados de
estado\footnote{La diferencia entre estados y otros tipos de predicados se
desarrolló en el \capref{8}.} que manifiestan una cualidad inherente y
caracterizadora del tema; así, por ejemplo, en \exref{ex:10-problema},
interpretamos que, si el problema se soluciona en dos minutos, se trata de un
problema \sig{fácil} o \sig{solucionable}. Dado que la propiedad referida se
presenta como atemporal, estas oraciones no se construyen en tiempos
perfectivos, que denotan que el evento efectivamente ocurrió, sino en tiempos
imperfectivos (presente o imperfecto), que indican que es posible que el evento
ocurra, de donde se desprende su valor modal\footnote{Sobre la lectura modal de
estas construcciones, v. \capref{14}.} \parencite{sanchezlopez2002}. Tanto los
tiempos imperfectivos como la presencia de adverbios o SSPP modales
(\forma{fácilmente, bien, con dificultad}) son inductores que promueven la
\termino{lectura caracterizadora} atemporal de estas construcciones.

\subsubsection{Construcciones cuasirreflejas pasivas (pasivas cuasirreflejas o pasivas con \textit{se})}\label{sec:10-pasivas}

Observemos las relaciones entre las construcciones de las columnas A, B y C.

\begin{center}\small
\begin{tabularx}{\linewidth}{@{}r L L L@{}}
  \toprule
  & A & B & C \\
  \midrule
  \refstepcounter{exx}(\theexx)\label{ex:10-traductor} & Un traductor novato tradujo el último libro
    premiado en el concurso Planeta. & El último libro premiado en el concurso
    Planeta fue traducido (por un traductor novato). & Se tradujo el último
    libro premiado en el concurso Planeta. \\
  \refstepcounter{exx}(\theexx) & Uno de los abogados filtró a la prensa las declaraciones del imputado. &
    Las declaraciones del imputado fueron filtradas a la prensa (por uno de los
    abogados). & Se filtraron a la prensa las declaraciones del imputado. \\
  \refstepcounter{exx}(\theexx) & Una empresa extranjera construirá una nueva represa en el sur del país. &
    Una nueva represa será construida en el sur del país (por una empresa
    extranjera). & Se construirá una nueva represa en el sur del país. \\
  \bottomrule
\end{tabularx}
\end{center}

Mientras las oraciones de la columna A son todas transitivas, ¿qué sucede con las
de B y C? En todos los casos el verbo sufre un proceso de intransitivización que
promueve el tema o paciente a la posición de sujeto y produce una
\termino{lectura de afectación}, porque el tema o paciente, que es el afectado
por el evento, pasa a ser más prominente. Sin embargo, debido a la
inacusativización de la estructura, el sujeto suele aparecer pospuesto, sobre
todo cuando es indeterminado. En cuanto al agente, en las oraciones de B puede
omitirse o aparecer marginalmente como complemento agente; en cambio, las de C
no admiten su presencia.\footnote{Aunque la presencia del agente es excepcional
con la pasiva con \forma{se}, en algunos textos jurídicos puede aparecer.}
Tradicionalmente, las oraciones de B se conocen como \termino{pasivas
perifrásticas}\footnote{Los conceptos de perífrasis verbal y de verbo auxiliar
se desarrollan en el \capref{11}.} porque el verbo adopta una forma compuesta
por el verbo auxiliar \forma{ser} (en el tiempo correspondiente de la voz activa
y concordado en número y persona con el sujeto tema) seguido del participio del
verbo principal (también concordado con el sujeto, en este caso, en género y
número). En las construcciones de C, que son \termino{pasivas cuasirreflejas},
el verbo está en forma activa y el clítico es un IPCR.

La eliminación sintáctica del agente no debe hacernos pensar que estas oraciones
no tienen un agente implícito que representa a \sig{quien realiza la acción}.
Una oración como \juicio{?}{\forma{La nueva represa se construirá sola}} resulta
extraña porque las pasivas, tanto perifrásticas como cuasirreflejas, involucran
eventos de acción transitivos que son realizados por un agente al que se puede
hacer referencia mediante un adverbio \exref{ex:10-deliberadamente} o una
construcción final \exref{ex:10-convenio}.

\begin{exe}
  \ex\label{ex:10-deliberadamente} Fue adelantada / Se adelantó la fecha de la
      reunión \textit{deliberadamente}. $\rightarrow$ Alguien la adelantó.
  \ex\label{ex:10-convenio} Se garantizó un nuevo convenio de precios \textit{para
      favorecer a los sectores de menores ingresos}. $\rightarrow$ Alguien lo
      garantizó.
\end{exe}

Sin embargo, no existe un paralelismo absoluto entre la pasiva perifrástica y la
pasiva con \forma{se}:

\begin{itemize}
  \item la primera admite todas las personas gramaticales, tanto singulares
        \exref{ex:10-estafado} como plurales \exref{ex:10-convocados}; la
        segunda está restringida a la tercera persona, tanto singular
        \exref{ex:10-prorrogo} como plural \exref{ex:10-prorrogaron}:
\end{itemize}

\begin{exe}
  \ex\label{ex:10-estafado} Fui estafado por uno de mis acreedores. /
      \juicio{*}{\forma{Se me estafó.}}\pendiente{En el español rioplatense,
      \textit{Se me estafó} es una impersonal aceptable}
  \ex\label{ex:10-convocados} Uds. serán convocados en el próximo llamado.
  \ex\label{ex:10-prorrogo} Se prorrogó la convocatoria.
  \ex\label{ex:10-prorrogaron} Se prorrogaron las fechas de examen.
\end{exe}

\begin{itemize}
  \item a diferencia de la pasiva perifrástica, que generalmente expresa eventos
        únicos \exref{ex:10-napoleon}, la pasiva con \forma{se} suele utilizarse
        en oraciones de sentido general \exref{ex:10-oportunidades}:
\end{itemize}

\begin{exe}
  \ex\label{ex:10-napoleon} Napoleón Bonaparte fue derrotado en 1815 en la
      batalla de Waterloo.
  \ex\label{ex:10-oportunidades} Se encuentran más oportunidades con una buena
      educación.
\end{exe}

Como la estructura pasiva exige ambos papeles ---tema y agente---, ni la forma
perifrástica ni la cuasirrefleja admiten verbos intransitivos, ya que los
inacusativos carecen de agente y los inergativos, de tema
\parencite{giammatteo2010}.

\begin{destacado}
\textbf{Para tener en cuenta.} \textcite{ditullio2014} plantea que, cuando el
sujeto de la pasiva cuasirrefleja es singular, se puede producir ambigüedad entre
la interpretación pasiva y la impersonal. En oraciones como \forma{Se convocó un
plebiscito} o \forma{Se requiere una auxiliar para tareas contables}, si se
considera que el SN acepta pronominalizarse, como sucede en el español
rioplatense (\forma{Se lo convocó} o \forma{Se la requiere}), la interpretación
es impersonal y el SN es COD. En cambio, si, dado que el SN concuerda con el
verbo en singular, se lo considera sujeto de la oración, la construcción tiene
interpretación pasiva.
\end{destacado}

\subsubsection{Construcciones cuasirreflejas impersonales}\label{sec:10-impersonales}

En español se consideran impersonales las oraciones que carecen de sujeto
explícito, ya sea porque el evento descripto excluye esa función, como en los
verbos meteorológicos (\forma{llover, tronar, granizar}), o porque el argumento
exigido carece de referencia definida, como en las \termino{construcciones
cuasirreflejas impersonales} con \forma{se}.

\begin{exe}
  \ex\label{ex:10-bebe} Aquí no \textit{se bebe}.
  \ex\label{ex:10-senado} En el Senado \textit{se debatió} durante largas horas.
  \ex\label{ex:10-sufre} \textit{Se sufre} si el amor no es correspondido.
  \ex\label{ex:10-esta} \textit{Se está} bien en este lugar.
  \ex\label{ex:10-cima} \textit{Se llega} a la cima muchas veces agotado.
  \ex\label{ex:10-descansa} Ya no \textit{se descansa} tan bien como antes en
      este lugar.
\end{exe}

¿Qué cambio se reconoce en la estructura argumental de estas oraciones? En todos
los ejemplos, los verbos suponen un sujeto ---agente en \exref{ex:10-bebe},
\exref{ex:10-senado}, \exref{ex:10-cima} y \exref{ex:10-descansa}; experimentante
en \exref{ex:10-sufre}, o paciente en \exref{ex:10-esta}---, aunque este no se
explicita y su referencia permanece indeterminada, ya que su mención se
considera innecesaria o irrelevante. Por tanto, el sujeto se puede representar
por un pronombre arbitrario.\footnote{La noción de \textit{pro} arbitrario se
introdujo en el \capref{6}, §\ref{sec:06-noverbales}.} La interpretación del
sujeto admite ya sea:

\begin{enumerate}[label=\alph*)]
  \item una \termino{lectura genérica}, parafraseable por \sig{todo el mundo /
        nadie / cualquiera}; por ejemplo, para \exref{ex:10-bebe}: \sig{Aquí
        nadie bebe}, o
  \item una \termino{lectura existencial}, parafraseable por \sig{alguno(s) /
        alguien / cualquiera}, como en \exref{ex:10-esta}: \sig{Cualquiera está
        bien en este lugar}.
\end{enumerate}

Los ejemplos muestran que no solo los verbos transitivos, \exref{ex:10-bebe} y
\exref{ex:10-senado}, admiten esta construcción, como sucede con la pasiva, sino
que prácticamente todos los tipos de verbos la aceptan, incluso los estativos,
\exref{ex:10-sufre} y \exref{ex:10-esta}, y los intransitivos inacusativos
\exref{ex:10-cima} e inergativos \exref{ex:10-descansa}.

\begin{recuadro}[Para ampliar]
Aunque algunos gramáticos han planteado que en las construcciones cuasirreflejas
impersonales el clítico \forma{se} representa el sujeto, esta hipótesis ha sido
descartada dentro del marco de la gramática generativa. Se postula, en cambio,
que el sujeto es un pronombre nulo que se especifica convencionalmente como
[\textit{pro}\etq{arb}] (con el subíndice \textit{arb}, \sig{arbitrario}) para
indicar que su referencia no es específica. Se considera que \forma{se} es el
elemento que desencadena la interpretación impersonal de la construcción, ya que
«absorbe» los rasgos definidos de persona y número de la flexión, que está
fijada en la forma no marcada de tercera persona singular, la más apta para
remitir a un sujeto inespecífico.\footnote{Para ampliar las nociones acerca de
la categoría gramatical de persona se puede consultar
\textcite{cunarro2018}.}
\end{recuadro}

Si bien en todas las construcciones cuasirreflejas impersonales ejemplificadas se
puede considerar que se ha producido un proceso de intransitivización, también
existen cuasirreflejas impersonales transitivas en las que, a diferencia de lo
que sucede en las cuasirreflejas medias, medio-pasivas y pasivas, el argumento
interno del verbo no es promovido a sujeto. La característica común de los
objetos de estas impersonales es que tienen el rasgo \sig{humano} o, al menos,
\sig{animado} y, para bloquear su interpretación como sujetos, reciben la marca
de función \forma{a} y pueden, además, ser pronominalizados.

\begin{exe}
  \ex Se arrestó a los sospechosos del triple crimen. $\rightarrow$ Se
      \textit{los} arrestó.
  \ex En la comisaría se metió a los imputados en una celda. $\rightarrow$ Se
      \textit{los} metió en una celda.
  \ex Al día siguiente, se convocó al abogado. $\rightarrow$ Se \textit{lo}
      convocó.
\end{exe}

\begin{recuadro}[Para ampliar: concordancias anómalas en las impersonales cuasirreflejas]
A pesar de que la cuasirrefleja impersonal está fijada en tercera persona
singular, en el uso surgen construcciones que «las gramáticas del español
deploran de forma unánime» \parencite[1689]{mendikoetxea1999}, con
concordancias «anómalas» desde el punto de vista normativo, como \forma{En el
ingreso al estadio se cachean a todos los que van entrando} o \forma{Se
encerraron a los culpables}. Pero ¿qué sucede en estas oraciones? Estos casos
pueden ponerse en línea con otros tan frecuentes como \forma{Hubieron fuertes
tormentas en el sur} o \forma{Habían muchas personas en la manifestación de
ayer}, en los que el hablante hace concordar un verbo impersonal como
\forma{haber} con su COD. Si bien para la normativa estas oraciones son
incorrectas, muchos hispanohablantes no podrían fácilmente detectar el error.
Según se plantea en \textcite{giammatteo2010}, en todas estas construcciones
está operando una tendencia regularizadora que lleva al hablante a «rellenar el
hueco estructural con lo que dentro de la oración mejor califica para sujeto»,
que es el único SN disponible: el COD. Volviendo al caso de las impersonales con
\forma{se} con objeto animado plural, un mecanismo interpretativo similar admite
la presencia de una doble señal: por un lado, la marca funcional \forma{a}, que
bloquea la interpretación agentiva del objeto y, al mismo tiempo, el plural en
el verbo, que señala que el paciente es reinterpretado por el hablante como
sujeto sintáctico de la oración.
\end{recuadro}

\begin{ejercicio}\label{ej:10-cuadro}
Complete el siguiente cuadro con las características comunes y diferenciales de
los cuatro subgrupos de construcciones cuasirreflejas que producen cambio de
estructura argumental, según se muestra en la primera columna.

{\small
\begin{tabularx}{\linewidth}{@{}L c c c c@{}}
  \toprule
  Características & Medias & Medio-pasivas & Pasivas & Impersonales \\
  \midrule
  Omisión del agente (o democión a función periférica) & + & & & \\
  Promoción del tema a sujeto de la oración & + & & & \\
  Tipo de lectura promovida & proceso & & & \\
  Persona de la concordancia & todas\textsuperscript{a} & & & \\
  Número de la concordancia & sg. y pl. & & & \\
  Presencia de inductores & $-$ & & & \\
  Naturaleza del predicado resultante & proceso & & & \\
  \bottomrule
\end{tabularx}}

\textsuperscript{a} \pendiente{El original dice «3.ª persona», aunque el texto
afirma que las medias son paradigmáticas}
\end{ejercicio}

\begin{ejercicio}\label{ej:10-meme}
Preste atención al siguiente meme y responda lo pedido en las consignas de
abajo.

\imagenpendiente{Meme con la construcción \textit{se veía}}

\begin{consignas}
  \item La construcción pronominal \forma{se veía}, que aparece en el meme,
        ofrece dos lecturas posibles. ¿Podría explicarlas y decir qué tipo de
        construcción pronominal es cada una de ellas?\\
        \textit{Lectura 1}: \forma{En el espejo se veía a sí mismo venir}.
        Construcción pronominal refleja pura.\\
        \textit{Lectura 2}: \completar[3]
  \item Luego explique la incidencia de ambas interpretaciones en el efecto
        humorístico del meme.
\end{consignas}
\end{ejercicio}

\begin{ejercicio}\label{ej:10-criolipolisis}
Analice las construcciones con \forma{se} de la siguiente publicidad y resuelva
lo pedido en las consignas, según se muestra en el ejemplo resuelto de abajo.

\begin{consignas}
  \item Discuta si es posible más de una lectura y aporte pruebas que lo
        fundamenten.
  \item ¿Qué funcionamiento tienen estas estructuras con \forma{se} en el texto?
\end{consignas}

\begin{quote}\small
\textbf{Criolipólisis}\\
Se ofrece un tratamiento para eliminar adiposidad localizada. Es no invasivo. Se
realiza una vez por mes y se requiere de 4 a 6 sesiones. Cada zona se trata
durante 15 minutos. Precio promocional 2x1 (la segunda persona se atiende sin
cargo).
\end{quote}

\begin{resuelto}
\forma{Se ofrece}: pasiva cuasirrefleja. No puede interpretarse como impersonal,
dado que acepta la versión en plural: \forma{Se ofrecen tratamientos}.
\end{resuelto}
\end{ejercicio}

\begin{ejercicio}\label{ej:10-shua}
Lea el texto «Ejemplares raros», de \textit{Botánica del caos}, de Ana María
Shua (2000). Observe las estructuras subrayadas y luego complete el cuadro de
abajo según se muestra en los casos resueltos.

% Texto literario protegido: se incorpora a mano desde el original.
\textoaincorporar{el microrrelato «Ejemplares raros», de A.~M.~Shua
(\textit{Botánica del caos}, 2000), con las estructuras subrayadas que figuran
en el cuadro}

{\small
\begin{tabularx}{\linewidth}{@{}>{\raggedright\arraybackslash}p{24mm}
    >{\raggedright\arraybackslash}p{13mm} L
    >{\raggedright\arraybackslash}p{15mm} >{\raggedright\arraybackslash}p{16mm}@{}}
  \toprule
  Construcción & Tipo de verbo & Cambio de la estructura argumental & Tipo de
    construcción pronominal & Tipo de lectura \\
  \midrule
  \forma{Se les frota el cuerpo} & acción & El agente es eliminado y el tema es
    promovido a sujeto. Acepta plural: \forma{Se les frotan los cuerpos}. &
    pasiva & afectación \\
  \forma{Se aplica (\ldots) una pomada antiséptica} & & & & \\
  \forma{Se las cubre} (a las llagas) & acción & El sujeto agente es eliminado,
    pero el tema no es promovido a sujeto y se mantiene en la función de COD. No
    acepta plural: \juicio{*}{\forma{Se las cubren a las llagas}}. & impersonal
    & existencial (\sig{alguien las cubre}) \\
  \forma{Se cierran} (las llagas) \forma{sin dejar marcas} & & & & \\
  \forma{las llagas se inflaman} & & & & \\
  \forma{Su movimiento se debe solo a la velocidad\ldots} & & & & \\
  \forma{Se las rocía} & & & & \\
  \forma{Se cubren} & & & & \\
  \forma{Las heridas se cierran} & & & & \\
  \bottomrule
\end{tabularx}}
\end{ejercicio}

\subsection{Construcciones cuasirreflejas que producen cambio en la estructura sintáctica}\label{sec:10-es}

En algunas construcciones pronominales cuasirreflejas de toda persona, el
pronombre indica un cambio en la estructura sintáctica que afecta a la
manifestación de la transitividad. En otras palabras, se cambia la relación
sintáctica que establece el verbo con su argumento interno, que pasa de ser COD a
complemento preposicional de régimen (CPR). Se distinguen dos subgrupos:

\begin{itemize}
  \item \termino{construcciones cuasirreflejas alternantes}, con verbos que
        permiten las dos construcciones, con COD y con CPR (\forma{olvidar
        (algo)} y \forma{olvidarse de (algo)});
  \item \termino{construcciones cuasirreflejas puras} o \termino{inherentes},
        con verbos que solo admiten la construcción pronominal (\forma{quejarse}
        / \juicio{*}{\forma{quejar}}; \forma{jactarse} /
        \juicio{*}{\forma{jactar}}).
\end{itemize}

\subsubsection{Construcciones cuasirreflejas alternantes}\label{sec:10-alternantes}

Observemos lo que sucede si comparamos las estructuras siguientes:

\begin{exe}
  \ex\label{ex:10-olvide}
    \begin{xlist}
      \ex Olvidé lo que quería decirte.
      \ex[*]{Olvidé de lo que quería decirte.}
      \ex Me olvidé de lo que quería decirte.
    \end{xlist}
  \ex\label{ex:10-ocupa}
    \begin{xlist}
      \ex Ocupa dos estantes de la librería.
      \ex Me ocupo de dos estantes de la librería.
    \end{xlist}
\end{exe}

Las construcciones cuasirreflejas que producen un cambio en la manifestación de
la transitividad tienen un comportamiento sintáctico homogéneo: la incorporación
del elemento pronominal siempre convierte el COD en un CPR. Sin embargo, desde la
perspectiva semántica, no todas las construcciones se comportan igual: algunas
mantienen su significado y otras lo alteran. Por ejemplo, en
\exref{ex:10-olvide}, el cambio sintáctico no modifica el significado del verbo,
mientras que, en \exref{ex:10-ocupa}, la alternancia no solo implica un cambio de
transitividad (de directa a indirecta), sino que también ocasiona un cambio en el
significado del verbo: en (a), \forma{ocupar} quiere decir \sig{llenar un
espacio} y, en (b), \forma{ocuparse de} significa \sig{asumir la responsabilidad
de un asunto}.\pendiente{El original remite a (88) y (89)}

\begin{ejercicio}\label{ej:10-alternantes}
Contraste las formas transitivas con las formas pronominales de los verbos
cuasirreflejos alternantes que figuran en la siguiente tabla y luego resuelva lo
pedido abajo.

{\small
\begin{tabularx}{\linewidth}{@{}l L l L@{}}
  \toprule
  Forma transitiva & Significado & Forma pronominal & Significado \\
  \midrule
  \forma{acordar}    & \sig{establecer un acuerdo} & \forma{acordarse de} &
    \sig{recordar} \\
  \forma{cruzar}     & & \forma{cruzarse con} & \\
  \forma{hallar}     & & \forma{hallarse en} & \\
  \forma{aprovechar} & & \forma{aprovecharse de} & \\
  \forma{empeñar}    & & \forma{empeñarse en} & \\
  \bottomrule
\end{tabularx}}

\begin{consignas}
  \item Según se muestra en el ejemplo resuelto, complete la tabla con el
        significado de cada uno de los verbos y luego redacte oraciones con cada
        uno de los significados de cada verbo.
  \item ¿Qué diferencia observa en el cambio de significado de los pares de 1
        frente a los de 2?
        \begin{enumerate}[label=\arabic*.]
          \item \forma{acordar / acordarse de; ocupar / ocuparse de; empeñar /
                empeñarse en}
          \item \forma{lamentar / lamentarse de; compadecer / compadecerse de;
                arriesgar / arriesgarse a}
        \end{enumerate}
\end{consignas}
\end{ejercicio}

\subsubsection{Construcciones cuasirreflejas puras o inherentes}\label{sec:10-inherentes}

Algunos verbos son inherentemente pronominales y no presentan alternancia. Solo
admiten la construcción pronominal, de modo que el clítico \forma{se} forma
parte del lexema. Dentro de estos verbos, algunos carecen de régimen verbal
(\forma{rebelarse, vanagloriarse, enfadarse, suicidarse, acurrucarse,
portarse}, etc.) y otros exigen un CPR (\forma{arrepentirse de, atreverse a,
dignarse a, quejarse de, abstenerse de, esforzarse en / por, jactarse de},
etc.).

\begin{exe}
  \ex
    \begin{xlist}
      \ex Se queja de todo permanentemente.
      \ex[*]{Queja de todo permanentemente.}
    \end{xlist}
  \ex
    \begin{xlist}
      \ex Estaba deprimido y se suicidó.
      \ex[*]{Estaba deprimido y suicidó.}
    \end{xlist}
\end{exe}

\begin{ejercicio}\label{ej:10-inherentes}
Analice las siguientes estructuras y determine qué cambios introduce el
pronombre, según se muestra en el ejemplo resuelto~1.

\begin{exe}
  \ex
    \begin{xlist}
      \ex Se atrevieron a decir la verdad.
      \ex[*]{Atrevieron a decir la verdad.}
    \end{xlist}
  \ex
    \begin{xlist}
      \ex La profesora despedía a sus alumnos.
      \ex La profesora se despedía de sus alumnos.
    \end{xlist}
  \ex
    \begin{xlist}
      \ex Cruzaron el bosque para llegar al lago.
      \ex Se cruzaron con ciervos de los pantanos.
    \end{xlist}
  \ex
    \begin{xlist}
      \ex Se dignaron venir por la mañana.
      \ex[*]{Dignaron venir por la mañana.}
    \end{xlist}
  \ex
    \begin{xlist}
      \ex Aseguró el clavo en la pared.
      \ex Se aseguró de la verdad antes de venir.
    \end{xlist}
\end{exe}

\begin{resuelto}
1. La agramaticalidad de 1b muestra que se trata de un caso de cuasirrefleja
pura o inherente: solo es posible la variante pronominal (1a), por lo que se
entiende que el pronombre forma parte del lexema.
\end{resuelto}
\end{ejercicio}

\subsection{Construcciones cuasirreflejas que producen cambio en la estructura eventiva}\label{sec:10-ee}

Para considerar los casos de este tercer grupo, podemos partir de los siguientes
ejemplos:

\begin{exe}
  \ex\label{ex:10-vino} ¿Quién (se) ha tomado todo el vino? (canción de La Mona
      Jiménez)
  \ex\label{ex:10-telon} (Se) cayó el telón.
  \ex\label{ex:10-duermo} Siempre (me) duermo temprano.
  \ex\label{ex:10-fui} (Me) fui hasta tu casa, pero no te encontré.
  \ex\label{ex:10-pipa} (Se) fumaron hasta la pipa de la paz.
\end{exe}

En todas estas construcciones se puede eliminar el pronombre y las oraciones
siguen siendo gramaticales. De ahí que se haya considerado que el pronombre es
un elemento espurio o expletivo, y que se haya buscado asimilar su función a la
del denominado \termino{dativo ético}.\footnote{Se denomina \textit{dativo ético}
al pronombre átono no reflexivo que se refiere al individuo implicado
afectivamente por la acción que denota el predicado
\parencite[102]{bosque2019}.} Sin embargo, mientras el dativo ético no concuerda
con el sujeto (\forma{Este chico no me come; No me vaya a llorar ahora};
ejemplos de \cite[102]{bosque2019}), el pronombre que aparece en las oraciones
de arriba concuerda en persona y número con el sujeto.

Por otra parte, desde el punto de vista semántico, ¿significan lo mismo las
versiones con y sin pronombre de las oraciones anteriores? Por ejemplo, cuando
decimos \forma{Cayó el telón} entendemos que el telón se bajó porque la obra
terminó; en cambio, \forma{Se cayó el telón} supone una caída accidental. Sin
embargo, en otros casos, las diferencias de significado entre la construcción con
y sin pronombre pueden ser más escurridizas.

El factor común en las oraciones de arriba es que la inclusión del pronombre no
elimina ni promueve argumentos, aunque puede incidir en su aceptabilidad. En
este tipo de construcciones pronominales, que pueden darse tanto con verbos
transitivos, \exref{ex:10-vino} y \exref{ex:10-pipa}, como con intransitivos,
\exref{ex:10-telon} a \exref{ex:10-fui}, la presencia del pronombre se asocia
con el aspecto léxico.\footnote{Las diferencias entre aspecto léxico y aspecto
gramatical y otras nociones relacionadas se tratan en el \capref{8}. Asimismo,
para ampliar estos conceptos, remitimos a \textcite{debin2018}.}

\begin{destacado}
\textbf{Para recordar.} El \termino{aspecto} es una categoría gramatical que
tiene que ver con la manera en que el evento se desarrolla y que tiene gran
incidencia sintáctica. Así, por un lado, circunscribe o delimita la estructura
interna del evento, lo que se conoce como \termino{aspecto cualitativo}, y, por
el otro, afecta su duración e intensidad, lo que corresponde al \termino{aspecto
cuantitativo}.
\end{destacado}

Pero ¿cómo inciden estas características del aspecto en las construcciones
pronominales? A partir de la distinción entre aspecto cualitativo, vinculado con
la delimitación, y aspecto cuantitativo, relacionado con la intensificación, se
reconocen dos subgrupos: construcciones cuasirreflejas de delimitación y
construcciones cuasirreflejas de intensificación.

\subsubsection{Construcciones cuasirreflejas de delimitación}\label{sec:10-delimitacion}

La diferencia de significado entre los verbos, tanto transitivos como
intransitivos, que se construyen con y sin pronombre que trataremos en este
apartado ha sido ampliamente debatida en la bibliografía, y el consenso al que
se ha llegado es que la variante pronominal restringe o delimita al verbo. Pero
¿en qué sentido? La incorporación del pronombre tiene valor aspectual y convierte
un verbo \termino{atélico} (no delimitado) en \termino{télico} (delimitado), con
un momento inicial y un estado resultante. Sin embargo, el efecto que esta
delimitación produce no es el mismo en todos los casos. Según vimos en los
ejemplos de arriba, estas construcciones pueden manifestar varios matices, todos
ellos vinculados con la delimitación, de modo que reconocemos dos subgrupos:
construcciones culminativas con verbos transitivos, como \exref{ex:10-vino} y
\exref{ex:10-pipa},\pendiente{El original incluye también (93)} y construcciones
ingresivas con verbos intransitivos, como \exref{ex:10-duermo} y
\exref{ex:10-fui}.

\paragraph{Construcciones culminativas con verbos transitivos.} Como ya hemos
señalado en el \capref{8},\pendiente{El original dice «capítulo 9»} el aspecto
es composicional porque no está dado por el verbo solamente, sino que múltiples
factores ---el tipo de sujeto, los complementos o algunos modificadores--- pueden
delimitar la realización del evento. Así, por ejemplo, a un verbo lo podemos
acotar con un COD determinado o cuantificado, como muestra el contraste entre
\exref{ex:10-manzanas-a} y \exref{ex:10-manzanas-b}, por un lado, y entre
\exref{ex:10-abuelo-a} y \exref{ex:10-abuelo-b}, por el otro:

\begin{exe}
  \ex
    \begin{xlist}
      \ex\label{ex:10-manzanas-a} Pedro come manzanas.
      \ex\label{ex:10-manzanas-b} Pedro comió una manzana / la manzana que estaba
          en la frutera / tres manzanas.
    \end{xlist}
  \ex
    \begin{xlist}
      \ex\label{ex:10-abuelo-a} El abuelo toma medicación para dormir.
      \ex\label{ex:10-abuelo-b} El abuelo tomó la / su medicación para dormir.
    \end{xlist}
\end{exe}

Solo en las construcciones en que el verbo aparece delimitado por el objeto
---\exref{ex:10-manzanas-b} y \exref{ex:10-abuelo-b}--- es posible el uso del
denominado \forma{se} \termino{aspectual}, que denota acciones «completas y
delimitadas» \parencite[111]{sanchezlopez2002}, lo que explica la
agramaticalidad de \exref{ex:10-secome-a} y \exref{ex:10-setoma-a} frente a
\exref{ex:10-secome-b} y \exref{ex:10-setoma-b}, que resultan
gramaticales.\pendiente{El original remite a «(108b)»}

\begin{exe}
  \ex
    \begin{xlist}
      \ex[*]{\label{ex:10-secome-a} Pedro se come manzanas.}
      \ex\label{ex:10-secome-b} Pedro se comió una manzana / la manzana que
          estaba en la frutera.
    \end{xlist}
  \ex
    \begin{xlist}
      \ex[*]{\label{ex:10-setoma-a} El abuelo se toma medicación para dormir.}
      \ex\label{ex:10-setoma-b} El abuelo se tomó su medicación / una pastilla
          para dormir.
    \end{xlist}
\end{exe}

La exigencia de un punto final que señale la culminación del evento hace que la
construcción pronominal no se emplee ni con actividades ni con estados.
Típicamente, estas construcciones se combinan con verbos de consumo
(\forma{comer, beber, fumar, gastar}), aunque también son compatibles con otros
verbos de realización que indiquen un estado final resultante. Asimismo, el
sujeto de estos verbos puede ser agente o experimentante (\forma{Facundo se sabe
el himno de memoria}), pero no paciente (\juicio{*}{\forma{Facundo se vive a la
vuelta de mi casa}}).

\paragraph{Construcciones ingresivas con verbos intransitivos.} Algunos verbos
intransitivos de movimiento (\forma{ir, venir, salir, llegar, correr, caer,
quedar, aparecer, volver}, etc.) y otros intransitivos de cambio de estado
(\forma{dormir, morir, despertar, estar}) presentan una construcción pronominal
(\forma{irse, venirse, salirse, llegarse, correrse, subirse, bajarse, caerse,
quedarse, aparecerse, volverse, dormirse, morirse, despertarse, estarse}, etc.).
Aunque todos estos verbos tienen en común el carácter intransitivo, las
construcciones pronominales que forman no muestran un comportamiento uniforme;
por eso vamos a ver ambos grupos separadamente.

\paragraph{Ingresivas con verbos intransitivos de movimiento.} Nuevamente vamos a
partir de la observación de algunos ejemplos contrastivos:

\begin{exe}
  \ex\label{ex:10-voy}
    \begin{xlist}
      \ex\label{ex:10-voy-a} Voy ([\etq{Origen} de la clase de Griego])
          [\etq{Meta} a la de Gramática].
      \ex[*]{\label{ex:10-voy-b} Voy [\etq{Origen} de la clase de Griego].}
      \ex Me voy [\etq{Origen} de la clase de Griego] ([\etq{Meta} a la de
          Gramática]).
    \end{xlist}
  \ex\label{ex:10-corrio}
    \begin{xlist}
      \ex\label{ex:10-corrio-a} Corrió ([\etq{Origen} de aquí]) [\etq{Meta} a la
          esquina].
      \ex[*]{\label{ex:10-corrio-b} Corrió [\etq{Origen} de aquí].}
      \ex Se corrió [\etq{Origen} de aquí].
      \ex[*]{Se corrió colina abajo.}
    \end{xlist}
\end{exe}

¿Qué efecto produce la inclusión del pronombre en las oraciones de arriba?
Prestemos atención a los papeles temáticos. Cuando estos verbos se construyen sin
pronombre, \exref{ex:10-voy-a} y \exref{ex:10-corrio-a}, el evento no se
encuentra delimitado y pueden aparecer ambos complementos, origen y meta, o solo
la meta, pero las oraciones resultan agramaticales cuando se construyen solo con
el argumento de origen, \exref{ex:10-voy-b} y \exref{ex:10-corrio-b}. La
inaceptabilidad de estas construcciones con un SP de origen muestra que, para
que el límite inicial pueda explicitarse, el verbo tiene que estar acompañado
por una forma pronominal, que «va a ser el elemento encargado de delimitar el
evento de \textit{ir} y hacerlo compatible con el SP delimitador»
\parencite[2986]{demiguel1999}. Para este subgrupo de verbos, el pronombre es
obligatorio cuando el verbo tiene lo que \textcite{demiguel1999} denomina
«presunción de origen» e indica abandono del lugar donde se encuentra el
hablante. Por tanto, al marcar el inicio del movimiento, estas construcciones se
consideran \termino{ingresivas} o \termino{incoativas}.

Según muestran los ejemplos anteriores, la presencia del pronombre no solo
influye en la aceptación de los complementos, sino que también incide en el
significado del verbo: mientras \forma{ir} es «moverse de un lugar hacia otro»
(\textit{DLE}), \forma{irse} significa «dejar un lugar (para ir a otro)»
\parencite[2986]{demiguel1999};\pendiente{El original cita aquí «De Miguel,
2009»} y frente a \forma{correr}, que significa \sig{ir deprisa},
\forma{correrse} (no registrado con esta acepción por el \textit{DLE}) es
\sig{salirse de un lugar} (Giammatteo, 2019).\pendiente{Giammatteo (2019) no
figura en la bibliografía}

No todos los verbos de movimiento presentan un claro contraste entre la forma
con y sin pronombre. Existen oposiciones más difusas, como \forma{Vino de la clase
de Griego a la de Gramática} y \forma{Se vino de la clase de Griego a la de
Gramática}, que suelen explicarse en función del grado de involucramiento del
sujeto, que sería mayor en la construcción pronominal. Incluso un mismo verbo
puede no presentar un comportamiento uniforme. Así, por ejemplo, son posibles
tanto \forma{Cayó de un quinto piso} como \forma{Se cayó de un quinto piso};
pero difiere la interpretación entre \forma{Cayó el pelotón de infantería}
(\sig{se rindió}) y \forma{Se cayó el pelotón de infantería (en la fosa)}
(\sig{se hundió}).

\begin{ejercicio}\label{ej:10-correr}
Observe las siguientes oraciones y resuelva lo solicitado en las consignas,
según se muestra en el ejemplo resuelto~1.

\begin{consignas}
  \item Explique los contrastes entre los casos a y b.
  \item Señale los papeles temáticos de todos los argumentos.
\end{consignas}

\begin{exe}
  \ex
    \begin{xlist}
      \ex ¡Corré!
      \ex ¡Correte de acá!
    \end{xlist}
  \ex
    \begin{xlist}
      \ex Cayó gente en la fiesta.
      \ex Se cayó gente en la fiesta.
    \end{xlist}
  \ex
    \begin{xlist}
      \ex La lluvia cae (al piso).
      \ex[*]{La lluvia se cae (al piso).}
    \end{xlist}
  \ex
    \begin{xlist}
      \ex Me cayó una maceta del balcón.
      \ex Se (me) cayó una maceta del balcón.
    \end{xlist}
  \ex
    \begin{xlist}
      \ex Se (me) cayó la silla.
      \ex Me caí de la silla.
    \end{xlist}
\end{exe}

\begin{resuelto}
1. a) En 1b, el pronombre \forma{te} tiene un valor aspectual y hace foco en el
inicio del evento. Se trata de una construcción ingresiva o incoativa que
produce cambio de significado: \forma{correr} es \sig{trasladarse con cierta
velocidad hacia algún otro lugar}, pero \forma{correrse} se interpreta como
\sig{abandonar un sitio para dirigirse a otro}. b) \forma{De acá} tiene el papel
temático de origen.
\end{resuelto}
\end{ejercicio}

\begin{ejercicio}\label{ej:10-cocina}
Observe las estructuras subrayadas del siguiente texto y determine los cambios
que produce la aparición del pronombre \forma{se}, según se muestra en el
ejemplo resuelto de abajo.

% Texto periodístico protegido: se incorpora a mano desde el original.
\textoaincorporar{el copete de la nota «Antes y después. Una cocina se actualiza
en blanco y en azul», con las estructuras subrayadas: \textit{se actualiza;
ampliarse; mudarse; quedarse}}

\begin{resuelto}
\forma{Una cocina se actualiza en blanco y en azul}. \forma{Se} produce cambios
en la estructura argumental del verbo: elimina al agente, que tiene
interpretación genérica (\forma{Cualquiera (sabe) que una cocina se actualiza en
blanco y azul}), y el tema es promovido a sujeto. Se pasa de una construcción
transitiva a una cuasirrefleja medio-pasiva con valor caracterizador de lo que el
texto plantea que \sig{debe ser una cocina actualizada}. El inductor del sentido
genérico de la construcción es el presente con valor modal deóntico (de
obligación) y, si bien no hay en el contexto un adverbio modal, por el sentido
puede reponerse \forma{siempre} $\rightarrow$ \forma{Una cocina debe ser
(siempre) actualizada en blanco y azul / Para ser actualizada, una cocina debe
estar en blanco y azul}.
\end{resuelto}
\end{ejercicio}

\paragraph{Ingresivas con verbos intransitivos de cambio de estado.} ¿Qué
diferencia hay entre \forma{dormir} y \forma{dormirse}? \forma{Dormir} es quizás
el caso más claro de los verbos de este grupo. Mientras que la forma no
pronominal refiere una actividad no delimitada, \forma{dormirse} es un verbo de
logro que marca el inicio del suceso, por lo que tiene valor ingresivo. Por
tanto, en uno y otro caso difiere la compatibilidad con modificadores
temporales.\footnote{La compatibilidad de las distintas clases de verbos con
modificadores temporales se presentó en el \capref{8},
§\ref{sec:08-mt}.}

\begin{exe}
  \ex
    \begin{xlist}
      \ex Estaba muy cansada y dormí profundamente toda la noche / *en un
          segundo.
      \ex Estaba muy cansada y me dormí en un segundo / *durante toda la noche.
    \end{xlist}
\end{exe}

Si bien existe cierto acuerdo en que la variante pronominal restringe los
valores del verbo y resulta aspectualmente marcada al denotar eventos
perfectivos, ¿cómo se aplica esto a verbos como \forma{morir} o
\forma{despertar}, que ya son télicos? Algo en común que tienen todos los verbos
intransitivos que admiten construcciones pronominales, tanto de movimiento como
de cambio de estado, es que solo las variantes pronominales permiten la
incorporación de un dativo no concordado con el sujeto, que señala al afectado
por el evento.

\begin{exe}
  \ex\label{ex:10-gato}
    \begin{xlist}
      \ex Murió el / mi gato. / \juicio{*}{\forma{Me murió el gato.}}
      \ex\label{ex:10-gato-b} Se murió el / mi gato. / Se me murió el gato.
    \end{xlist}
  \ex\label{ex:10-lapiz}
    \begin{xlist}
      \ex Cayó el lápiz. / \juicio{*}{\forma{Me cayó el lápiz.}}
      \ex\label{ex:10-lapiz-b} Se cayó el lápiz. / Se me cayó el lápiz.
    \end{xlist}
\end{exe}

Varios de los verbos pronominales inacusativos admiten estas
\termino{construcciones doblemente pronominales} o \termino{bipronominales}, que
incorporan un dativo no concordado como segunda forma pronominal, cuya presencia
altera el orden de constituyentes porque el dativo se coloca entre el pronombre
concordado y el verbo \exref{ex:10-anteojos}:

\begin{exe}
  \ex\label{ex:10-anteojos}
      \begin{tabular}[t]{@{}llll@{}}
        Se & me & cayeron & los anteojos. \\
        \footnotesize clít. concordado & \footnotesize + clít. dat. no
          concordado & \footnotesize + V & \footnotesize + sujeto \\
      \end{tabular}
\end{exe}

En la forma no marcada de las construcciones doblemente pronominales, el sujeto
se pospone y el dativo, que representa al afectado por el evento, puede estar
reduplicado por un SP que suele colocarse al inicio:

\begin{exe}
  \ex A Juan se le cayó el cuaderno.
\end{exe}

\begin{recuadro}[Para ampliar: el dativo afectado]
En estas construcciones el dativo se considera el argumento más prominente, ya
que refiere a una entidad animada, mientras que el sujeto pospuesto suele ser no
animado. Para \textcite[9]{maldonado1994}, este \termino{dativo de afectación}
es «una construcción intermedia cuyo análisis ha sido problemático a lo largo de
la historia lingüística del español». Según \textcite[103]{sanchezlopez2002},
«aunque dativos, estos argumentos no pueden considerarse benefactivos, sino que
representan un argumento que es en cierto modo responsable del proceso sin
intervenir en él directamente». Otros autores
\parencite{masullo1992,fernandezsoriano1999,mendikoetxea2011} los consideran
sujetos en dativo, ya que pueden ocupar la posición canónica inicial y hacer las
veces de «sujeto psicológico» de la oración.\pendiente{El original cita
«Fernández Soriano, 1990»} El dativo afectado se combina con distintos tipos de
verbos que tienen en común el hecho de no poseer un sujeto que cause el evento,
como sucede con:

\begin{itemize}
  \item las formas pronominales de los verbos de cambio de estado de causa
        externa (\forma{quemar, estropear, romper}), con los que el dativo
        representa al responsable indirecto del evento;
        \begin{exe}
          \ex El jarrón se me rompió.
        \end{exe}
  \item las formas pronominales de verbos intransitivos inacusativos
        (\forma{morirse} \exref{ex:10-gato-b}, \forma{caerse}
        \exref{ex:10-lapiz-b}, \forma{irse, salirse}) con «presunción de
        origen». Como muestra \exref{ex:10-diente}, la incorporación del
        pronombre puede implicar un cambio de significado en el
        verbo;\pendiente{El original remite a (107)}
        \begin{exe}
          \ex\label{ex:10-diente}
            \begin{xlist}
              \ex Le salió un diente nuevo al bebé. (\sig{apareció})
              \ex Se le salió un diente al abuelo. (\sig{se desprendió de la
                  encía})
            \end{xlist}
        \end{exe}
  \item predicaciones impersonales (\forma{llover, anochecer}).
        \begin{exe}
          \ex Nos llovió toda la semana y tuvimos que volver de las vacaciones.
        \end{exe}
\end{itemize}

Con los inacusativos pronominales y los impersonales meteorológicos, el dativo
representa al experimentante afectado psicológicamente por el suceso.
\textcite{company2017} plantea que las construcciones con dativo afectado
constituyen una estrategia de subjetivización del hablante para interpretar los
sucesos desde su óptica personal y de manera afectiva. Así, aunque con
\forma{quemar} la construcción agentiva es perfectamente gramatical, una
oración como \exref{ex:10-asado-a} sería poco probablemente emitida por un
hablante. Por lo general, se prefiere la formulación anticausativa
\exref{ex:10-asado-b}, que presenta al evento como no causado y omite al
responsable. La aparición del dativo \exref{ex:10-asado-c} posibilita que el
causante se presente como un afectado por el hecho, en una lectura de proceso
accidental, en la que el sujeto se muestra como responsable indirecto o
involuntario del suceso, aunque involucrado en él.

\begin{exe}
  \ex
    \begin{xlist}
      \ex\label{ex:10-asado-a} Yo quemé el asado. ⟨agente, tema⟩
      \ex\label{ex:10-asado-b} El asado se quemó. ⟨tema⟩
      \ex\label{ex:10-asado-c} [Me distraje y] el asado se me quemó. ⟨tema,
          experimentante⟩
    \end{xlist}
\end{exe}

Los distintos casos que hemos presentado permiten ver el complejo estatus de
estos dativos, tradicionalmente considerados superfluos, pero que, sin embargo,
son determinantes a la hora de fijar la interpretación de la oración.
\end{recuadro}

Asimismo, como ya vimos con algunos verbos intransitivos de movimiento, también
con los de cambio la diferencia entre la construcción con y sin pronombre suele
estar dada porque la primera hace hincapié en el esfuerzo o la intencionalidad
del sujeto en la realización del evento. Así, en algunos casos se otorga
agentividad a sujetos que en las construcciones no pronominales se considerarían
pacientes, tal como muestran los contrastes siguientes:

\begin{exe}
  \ex
    \begin{xlist}
      \ex Vino a mi casa el domingo.
      \ex Se vino a mi casa desde José C. Paz en tren.
    \end{xlist}
  \ex
    \begin{xlist}
      \ex Apareció la niña que estaban buscando hacía varios días
          (casualmente).
      \ex Se apareció en la fiesta, aunque no estaba invitado (deliberadamente).
    \end{xlist}
  \ex
    \begin{xlist}
      \ex Trajo una torta para festejar.
      \ex Se trajo una torta para festejar.
    \end{xlist}
\end{exe}

\subsubsection{Construcciones cuasirreflejas de intensificación}\label{sec:10-intensificacion}

Estas construcciones se dan tanto con verbos transitivos como con intransitivos
inergativos y, según \textcite{ditullio2014}, son propias de la variedad
rioplatense. La incorporación del pronombre, apoyada por la presencia de algún
elemento intensificador, señala «[\ldots] un grado extremo en que se realiza una
actividad y enfatiza, así, la capacidad o habilidad del agente para llevarla a
cabo» \parencite[181]{ditullio2014} \exref{ex:10-traviata}, aunque el sujeto
también puede ser un experimentante, \exref{ex:10-floreros} y
\exref{ex:10-sabian}.

\begin{exe}
  \ex\label{ex:10-traviata} Se cantó todo / la vida / hasta \textit{La
      traviata}.
  \ex\label{ex:10-floreros} Después del viaje teníamos tanta sed que nos tomamos
      todo / hasta el agua de los floreros.
  \ex\label{ex:10-sabian} Se sabían todo / vida y milagros del cantante.
\end{exe}

\begin{ejercicio}\label{ej:10-intensificacion}
Analice los cambios producidos en las siguientes estructuras a partir de la
incorporación del pronombre, según se muestra en el ejemplo resuelto~1.

\begin{exe}
  \ex a) Tomó café. / b) *Se tomó café. / c) Se tomó un café. / d) Se tomó el
      café de su vida.
  \ex Se sabe todo, pero no lo cuenta a nadie.
  \ex a) Se comió toda la carne. / b) Se comió hasta las migas.
  \ex a) Fuma desde chico. / b) ¡Hoy se fumó la vida! / c) Se fumó varios
      cigarrillos durante el examen.
  \ex a) Salió con sus hijos al parque. / b) Se salió el botón de la camisa. /
      c) Me salí de la dieta.
  \ex Se iban a la iglesia a pedir perdón. (Videograf del canal televisivo
      \textit{Crónica})
\end{exe}

\begin{resuelto}
1. a) Oración transitiva. b) Oración agramatical: el COD es un nombre de materia
no determinado ni cuantificado que no delimita el evento, por lo que no es
posible la incorporación del \forma{se} aspectual delimitador. c) Construcción
cuasirrefleja pronominal culminativa de delimitación: el COD está delimitado y
la incorporación del \forma{se} aspectual es gramatical. d) Construcción
cuasirrefleja pronominal de intensificación que se interpreta como \sig{el mejor
café que tomó en su vida}, con valor superlativo.
\end{resuelto}
\end{ejercicio}

\begin{ejercicio}\label{ej:10-diferencias}
Identifique en cada caso la construcción pronominal involucrada y explique las
diferencias entre las versiones a y b de cada ejemplo, según se muestra en el
ejemplo resuelto~1.

\begin{exe}
  \ex
    \begin{xlist}
      \ex Se buscan especialistas en manejo de herramientas informáticas.
      \ex Se busca a los especialistas interesados en tomar becas.
    \end{xlist}
  \ex
    \begin{xlist}
      \ex Para las fiestas se come mucho.
      \ex Se comieron la entrada, el plato principal y el postre.
    \end{xlist}
  \ex
    \begin{xlist}
      \ex Este tema se aprende con facilidad.
      \ex Se aprendió de memoria la lección.
    \end{xlist}
  \ex
    \begin{xlist}
      \ex Se están enviando turnos de vacunación.
      \ex Los alumnos se enviaron mensajes anónimos.
    \end{xlist}
\end{exe}

\begin{resuelto}
1. La incorporación del pronombre \forma{se} en ambas construcciones provoca un
cambio en la estructura argumental: el agente es eliminado. Las diferencias
estriban en que en a el tema (\forma{especialistas\ldots}) es promovido a
sujeto, lo que permite una lectura pasiva; en b, en cambio, la marca de función
\forma{a} impide la promoción del tema (\forma{los especialistas
interesados\ldots}) a la función de sujeto, lo que produce una lectura
impersonal.
\end{resuelto}
\end{ejercicio}

\begin{ejercicio}[Actividad de integración]\label{ej:10-akasha}
Observe las estructuras subrayadas en el siguiente fragmento del cuento
«Akasha», de S.~Fernández (2019), y resuelva las consignas que se presentan a
continuación del texto, según se muestra en los ejemplos resueltos de abajo.

% Texto literario protegido: se incorpora a mano desde el original.
\textoaincorporar{el epígrafe y el fragmento de «Akasha», de S.~Fernández
(2019), desde «Cierro mis ojos para descansar» hasta «\ldots{}esta nueva
sensación de desamparo», con las estructuras pronominales subrayadas}

\begin{consignas}
  \item En cada uno de los casos, identifique si la referencia del pronombre
        coincide con la del sujeto y determine si la construcción es pronominal
        o no.
  \item En las construcciones pronominales, evalúe si el pronombre personal
        involucrado manifiesta un argumento del verbo al que acompaña o no:
        \begin{enumerate}[label=\roman*)]
          \item en caso afirmativo, determine el papel temático y la función
                sintáctica del clítico y proporcione pruebas para identificar de
                qué tipo de pronombre reflejo se trata (reflejo puro, recíproco
                o factitivo);
          \item en caso negativo, cuando el pronombre involucrado no manifieste
                ningún argumento del verbo, indique qué tipo de cambio introduce
                (en la estructura sintáctica, en la argumental o en la
                eventiva), de qué tipo de construcción cuasirrefleja se trata y
                aporte pruebas que justifiquen su decisión.
        \end{enumerate}
\end{consignas}

\begin{resuelto}[Ejemplos resueltos]
\forma{(El aire) me molesta}. a) La referencia del pronombre no coincide con la
del sujeto y la construcción no es pronominal. \forma{(El poder espiritual de
carácter omnipresente) se encuentra\ldots}: a) la referencia del pronombre
coincide con la del sujeto y la construcción es pronominal. b) El pronombre no
manifiesta ningún argumento del verbo. ii) El pronombre produce cambio en la
estructura argumental y la construcción es cuasirrefleja pasiva porque admite
tercera persona singular y plural: \forma{Los poderes espirituales de carácter
omnipresente se encuentran}.\pendiente{El original dice aquí «estructura
sintáctica»; el solucionario, «estructura argumental»}
\end{resuelto}
\end{ejercicio}

\begin{recuadro}[Las construcciones pronominales en el texto]
En este fragmento del cuento se observa que la narradora construye una escena en
la que los eventos que se suceden «le suceden» al personaje. Aun en los casos en
los que a la protagonista le corresponde el papel de agente, parece actuar
dejándose llevar por las circunstancias (\forma{cierro (mis ojos), vuelo,
extiendo (mis brazos), rozo (algo), vuelvo (al tren), digo, viajo (lo hago por
primera vez)}) o a la expectativa de lo que le viene del exterior (\forma{miro,
escucho}). Para lograr el clima de despersonalización que la lleva de la
ensoñación y la sensación de libertad inicial a la soledad y el desamparo final
---un pasaje de la vida a la muerte, como prefigura el epígrafe---, la
protagonista cada vez asume una actitud más pasiva, que deja ver un personaje
inseguro, al cual «las cosas le ocurren». Este efecto se logra porque predominan
los casos en que ella no es el agente que lleva adelante la acción, sino que
aparece ya sea como objeto paciente de verbos transitivos (\forma{(el aire)
\textbf{me} rodea, \textbf{me} eleva; (la puerta) \textbf{me} saca del ensueño;
(mi acompañante) \textbf{me} mira; \textbf{me} despierta (una inmensa
curiosidad); \textbf{me} hace sentir; \textbf{me} pregunto}), como sujeto
paciente de verbos de estado (\forma{estoy sentada, estoy en un tren, no estoy
tan sola}) o como experimentante de verbos psicológicos (\forma{siento,
\textbf{me} molesta, creo, no sé, vuelvo a pensar}). Otro recurso de
despersonalización empleado en el fragmento es el uso de construcciones
pronominales medias, que favorecen la presencia de sujetos no agentivos
(\forma{se deslizan (mis dedos), se borra (mi sonrisa), me levanto}). Junto con
las construcciones verbales, refuerzan la actitud de inseguridad de la
protagonista las preguntas y los índices de duda y negación que indican
desconocimiento de la situación (\forma{¿Cuánto hace que estoy sentada? No lo
sé. ¿Cómo lo sé? No tengo idea. ¿A dónde voy? ¿Acaso importa?}, \forma{quizás
sea quien respondió, no estoy tan sola, quizás, a dónde voy, a dónde vamos}).
Asimismo, también los elementos léxicos acompañan esta presentación y la
protagonista pasa de una \forma{ensoñación} en la que se siente \forma{libre} a
un estado en que \forma{se le borra la sonrisa} y \forma{se siente incómoda, con
malestar}, hasta finalmente sentir \forma{la soledad y el desamparo}.
\end{recuadro}

\begin{recuadro}[Reflexión final]
El recorrido realizado en este capítulo nos ha mostrado que el campo de las
construcciones pronominales es sumamente amplio y que las distintas subclases no
siempre presentan bordes nítidos, lo que las convierte en uno de los temas más
complejos de la gramática del español, aunque también, según creemos, en uno de
los más interesantes. En el capítulo hemos distinguido dos grupos principales:
las reflejas, en las que el pronombre tiene función sintáctica de COD o de
COI,\pendiente{El original dice «de COD o de COD»} y las cuasirreflejas, que se
parecen a las primeras en la forma, pero en las que el pronombre no cumple
ninguna función sintáctica. Si bien en las reflejas, dentro de las que hemos
distinguido entre las reflexivas (puras, accidentales y factitivas) y las
recíprocas, se da un paso hacia la intransitivización, en tanto el agente y el
paciente coinciden, es en las cuasirreflejas en las que se advierte este proceso
como completado, debido a que se produce el borrado o la democión del agente
dentro de la estructura oracional. Con fines didácticos hemos adoptado para estas
construcciones cuasirreflejas una triple distinción, que tiene en cuenta si la
introducción del pronombre afecta a la estructura argumental (cuasirreflejas
medias, medio-pasivas, pasivas e impersonales), a la estructura sintáctica
(cuasirreflejas alternantes e inherentes) o a la estructura eventiva
(cuasirreflejas de delimitación y de intensificación).
\end{recuadro}

Por último, a modo de recapitulación, en la tabla~\ref{tab:10-sintesis} ofrecemos
una síntesis de las distinciones reconocidas en el capítulo.

{\footnotesize\setlength{\tabcolsep}{3pt}
\begin{longtable}{@{}>{\raggedright\arraybackslash}p{30mm}
    >{\raggedright\arraybackslash}p{38mm} >{\raggedright\arraybackslash}p{52mm}@{}}
  \caption{Sistematización de construcciones pronominales}\label{tab:10-sintesis}\\
  \toprule
  Grupo & Subtipo & Ejemplo \\
  \midrule
  \endfirsthead
  \toprule
  Grupo & Subtipo & Ejemplo \\
  \midrule
  \endhead
  \multicolumn{3}{@{}l}{\textbf{Reflejas}: el pronombre es COD o COI} \\
  \addlinespace
  Reflexivas (Pr Refl) & reflexivas puras & \textit{Juan se vio (a sí mismo).} \\
  & reflexivas accidentales & \textit{Juan se rompió el dedo
    (accidentalmente).} \\
  & reflexivas factitivas & \textit{Juan se cortó el pelo en la peluquería} (=
    \textit{Se hizo cortar el pelo en la peluquería}). \\
  Recíprocas (Pr Recipr) & --- & \textit{María y Pedro se miraron (el uno al
    otro / recíprocamente).} \\
  \midrule
  \multicolumn{3}{@{}p{120mm}}{\textbf{Cuasirreflejas}: el pronombre no tiene
    función sintáctica (es IPCR); producen cambios en la oración} \\
  \addlinespace
  Cambio en la estructura sintáctica & alternantes (con cambio de significado o
    sin él) & \textit{Olvidé eso / Me olvidé de eso. Ocupa dos estantes / Se
    ocupa de sus hijos.} \\
  & inherentes & \textit{Juan se arrepintió de lo que dijo.} \\
  \addlinespace
  Cambio en la estructura argumental: de toda persona & medias: verbos de cambio
    de estado físico (\textit{romper / romperse}), de cambio de estado mental
    (\textit{molestar / molestarse}) o de cambio de posición (\textit{sentar /
    sentarse}) & \textit{La puerta se cerró (sola). María se divirtió. El bebé ya
    se sienta (solo).} \\
  Cambio en la estructura argumental: de tercera persona & medio-pasivas &
    \textit{Esta carne se corta fácilmente.} \\
  & pasivas (el verbo concuerda con el sujeto en 3.ª persona, sg. o pl.) &
    \textit{Se tomó una medida polémica / Se tomaron medidas polémicas.} \\
  & impersonales (el verbo está fijado en 3.ª persona sg.) & \textit{Se habla
    mucho y se trabaja poco.} \\
  \addlinespace
  Cambio en la estructura eventiva o aspectual & de delimitación & \textit{Se
    comió toda la torta. Se fue a su casa. Se durmió temprano.} \\
  & de intensificación & \textit{Se comió hasta las migas.} \\
  \bottomrule
\end{longtable}}

% --- Secciones finales -------------------------------------------------
\lecturasrecomendadas

\solucionario

\begin{solucion}{\ref{ej:10-caso}}
2. \forma{Lo}: caso acusativo, función sintáctica COD. 3. \forma{Me}: caso
dativo, función sintáctica COI. 4. \forma{Me}: caso dativo, función sintáctica
COI. 5. \forma{Te}: caso dativo, función sintáctica COI; \forma{lo}: caso
acusativo, función sintáctica COD; \forma{me}: caso acusativo, función
sintáctica COD.
\end{solucion}

\begin{solucion}{\ref{ej:10-funcion}}
2. \forma{Le}: función sintáctica COI, papel temático meta. 3. \forma{Lo}:
función sintáctica COD, papel temático tema. 4. \forma{Le}: función sintáctica
COI, papel temático meta (\forma{le} duplica al COI \forma{a su hermano}). 5.
\forma{Se}: función sintáctica COI, papel temático meta. A diferencia de 4, aquí
aparecen ambos objetos pronominalizados, por lo que en lugar de \forma{le}
aparece el alomorfo \forma{se}. \forma{Lo}: función sintáctica COD, papel
temático tema.
\end{solucion}

\begin{solucion}{\ref{ej:10-sustitucion}}
2. a) Sustitución: \forma{Le arreglé el pelo frente al espejo}. b) El pronombre
sustituto \forma{le} manifiesta caso dativo; por lo tanto, la función sintáctica
de \forma{me} es COI. 3. a) Sustitución: \forma{Lo vistió de gala}. b) El
pronombre sustituto \forma{lo} manifiesta caso acusativo; por lo tanto, la
función sintáctica de \forma{se} es COD. 4. a) Sustitución: \forma{Lo ayudaron
con la tarea}. b) El pronombre sustituto \forma{lo} manifiesta caso acusativo;
por lo tanto, la función sintáctica de \forma{te} es COD. 5. a) Sustitución:
\forma{Les dieron la vacuna para viajar}. b) El pronombre sustituto \forma{les}
manifiesta caso dativo; por lo tanto, la función sintáctica de \forma{nos} es
COI.
\end{solucion}

\begin{solucion}{\ref{ej:10-papeles}}
2. \forma{Te}: paciente, COD; \forma{nos}: beneficiario, COI. 3. \forma{Se}:
beneficiario, COI; \forma{lo}: tema, COD. 4. \forma{Me}: beneficiario, COI;
\forma{la}: tema, COD. 5. \forma{Se}: beneficiario, COI; \forma{las}: tema,
COD.
\end{solucion}

\begin{solucion}{\ref{ej:10-factitivas}}
Mi primo \forma{se} hizo una hermosa casa: agente: la empresa o persona que
construye; causante instigador: el sujeto, \forma{mi primo}. El pronombre
\forma{se} es COI. \forma{Me} hice un piercing: agente: la persona que colocó el
piercing; causante instigador: el sujeto tácito \forma{yo}. El pronombre
\forma{me} es COI. Lorena \forma{se} vacunó: agente: la persona o enfermera que
colocó la vacuna; causante instigador: \forma{Lorena}. El pronombre \forma{se}
es COI.\pendiente{El original intercambia «agente» y «causante instigador» y
analiza \textit{se} como COI en \textit{se vacunó} (es COD)}
\end{solucion}

\begin{solucion}{\ref{ej:10-ambiguas}}
Respuestas posibles. a) 2. Reflexiva: \forma{Cada uno de ellos se pasó la
tintura por el propio pelo}. Recíproca: \forma{Ana le pasó la tintura a su
hermano mientras él le hacía algunos retoques de tintura a ella}. 3. Refleja
pura: \forma{Cada uno se copió usando su propio celular}. Recíproca: \forma{Cada
uno se copió del otro / Se copiaron mutuamente}. b) En 2 también es posible la
interpretación como refleja factitiva: \forma{Ana y su hermano se hicieron teñir
el pelo en la peluquería de la esquina de su casa}.
\end{solucion}

\begin{solucion}{\ref{ej:10-tiporefleja}}
2. Reflexiva factitiva. 3. Reflexiva pura / factitiva. 4. Reflexiva pura. Esta
oración es ambigua entre una interpretación accidental y la reflexiva pura, pero
se desambigua por conocimiento de mundo, ya que, según se sabe, el famoso pintor
se automutiló. 5. Refleja recíproca.
\end{solucion}

\begin{solucion}{\ref{ej:10-danna}}
\forma{Se toca el cabello}: construcción pronominal refleja. El pronombre
\forma{se} es correferencial con el sujeto. Admite el refuerzo \forma{a sí
misma}. \forma{Le arregle las pestañas y las cejas}: no es construcción
pronominal. El pronombre \forma{le} no es correferencial con el sujeto.
\forma{Me estoy agarrando el pelo}: construcción pronominal refleja. El
pronombre \forma{me} es correferencial con el sujeto. Admite el refuerzo
\forma{a mí misma}. \forma{Me vean bien el pelo}: no es construcción pronominal.
El pronombre \forma{me} no es correferencial con el sujeto. \forma{Me las tengo
que hacer yo}: construcción pronominal refleja. El pronombre \forma{me} es
correferencial con el sujeto. Admite el refuerzo \forma{a mí misma}. \forma{Me
las haga}: no es construcción pronominal. El pronombre \forma{me} no es
correferencial con el sujeto. \forma{Se comió todas las uñas de la mano
derecha}: construcción pronominal refleja. El pronombre \forma{se} es
correferencial con el sujeto. Admite el refuerzo \forma{a sí misma}.
\end{solucion}

\begin{solucion}{\ref{ej:10-levantar}}
En 2a, el verbo \forma{extender} se comporta como transitivo: se construye con
sujeto agente (\forma{el nuevo gobierno}) y COD tema (\forma{los límites}), y
acepta la paráfrasis con \forma{hacer}: \forma{¿Qué hizo el gobierno?} En 2b
aparece la variante pronominal de \forma{extender}, referida a un suceso que
afecta a un tema. Se puede añadir el refuerzo \forma{solo} o \forma{por sí
solo}.
\end{solucion}

\begin{solucion}{\ref{ej:10-cambio}}
2. Verbo de cambio de posición. 2a: construcción transitiva con agente (tácito)
y tema (\forma{las cortinas}); 2b: construcción cuasirrefleja con sujeto
(tácito) simultáneamente agente y tema. 3. Verbo de cambio de posición. 3a:
construcción transitiva con agente (tácito) y paciente (\forma{los niños}); 3b:
construcción cuasirrefleja con sujeto (tácito) simultáneamente agente y tema.
\end{solucion}

\begin{solucion}{\ref{ej:10-cuadro}}
{\small
\begin{tabularx}{\linewidth}{@{}L c >{\centering\arraybackslash}p{20mm} c
    >{\centering\arraybackslash}p{18mm}@{}}
  \toprule
  Características & Medias & Medio-pasivas & Pasivas & Impersonales \\
  \midrule
  Omisión del agente (o democión a función periférica) & + & + & + & + \\
  Promoción del tema a sujeto de la oración & + & + & + & $-$ \\
  Tipo de lectura promovida & proceso & caracterizadora & afectación & genérica
    o existencial \\
  Persona de la concordancia & todas & 3.ª & 3.ª & 3.ª \\
  Número de la concordancia & sg. y pl. & sg. y pl. & sg. y pl. & sg. \\
  Presencia de inductores & $-$ & tiempos imperfectivos y adverbios modales &
    $-$ & $-$ \\
  Naturaleza del predicado resultante & proceso & estado & acción & sin
    restricciones \\
  \bottomrule
\end{tabularx}}
\end{solucion}

\begin{solucion}{\ref{ej:10-meme}}
a) Lectura 2: \forma{Se veía (en el espejo) esto: que él venía}. Construcción
pronominal cuasirrefleja pasiva. b) El efecto humorístico surge del contraste
entre ambas interpretaciones: en la lectura 1 se da a entender que el hombre
muere al chocar (él mismo) contra el espejo; en la 2, lo que se plantea es que
la muerte era previsible (= \forma{se (la) veía venir}).
\end{solucion}

\begin{solucion}{\ref{ej:10-criolipolisis}}
\forma{Se realiza}, \forma{se trata} y \forma{se atiende} son pasivas
cuasirreflejas. No pueden interpretarse como impersonales porque pueden
pluralizarse: \forma{Se realizan tratamientos; Todas las zonas se tratan durante
15 minutos; Las dos personas se atienden sin cargo}. \forma{Se requiere} es una
cuasirrefleja impersonal. No puede interpretarse como pasiva porque no acepta el
plural: \juicio{*}{\forma{Se requieren de 4 a 6 sesiones}}. Función textual de
estas construcciones: tanto las cuasirreflejas pasivas como las impersonales, al
promover el «borrado» del agente, resultan aptas para textos en los que cobra
mayor relevancia el tema que el agente.
\end{solucion}

\begin{solucion}{\ref{ej:10-shua}}
{\footnotesize\setlength{\tabcolsep}{3pt}
\begin{longtable}{@{}>{\raggedright\arraybackslash}p{24mm}
    >{\raggedright\arraybackslash}p{16mm} >{\raggedright\arraybackslash}p{44mm}
    >{\raggedright\arraybackslash}p{14mm} >{\raggedright\arraybackslash}p{18mm}@{}}
  \toprule
  Construcción & Tipo de verbo & Cambio de la EA & Tipo de construcción
    pronominal & Tipo de lectura \\
  \midrule
  \endhead
  \forma{Se aplica (\ldots) una pomada antiséptica} & acción & El agente es
    eliminado y el tema es promovido a sujeto. Acepta plural: \forma{Se aplican
    pomadas}. & pasiva & afectación \\
  \forma{Se cierran (las llagas) sin dejar marcas} & cambio de estado físico & El
    agente es eliminado y el tema es promovido a sujeto. & media & proceso \\
  \forma{Las llagas se inflaman} & cambio de estado físico & El agente es
    eliminado y el tema es promovido a sujeto. & media & proceso \\
  \forma{Su movimiento se debe solo a la velocidad (\ldots)} & acción & El agente
    es demovido a ACC y el tema es promovido a sujeto. Acepta plural: \forma{Sus
    movimientos se deben a (\ldots)}. & pasiva & afectación \\
  \forma{Se las rocía} & acción & El agente es eliminado, pero el tema no es
    promovido a sujeto y se mantiene en la función de COD. No acepta plural:
    \juicio{*}{\forma{Se las rocían (a las raicillas)}}. & impersonal &
    existencial (\sig{alguien las rocía}) \\
  \forma{Se cubren} & acción & El agente es eliminado y el tema es promovido a
    sujeto. Acepta singular: \forma{Se cubre}. & pasiva & afectación \\
  \forma{Las heridas se cierran} & cambio de estado físico & El agente es
    eliminado y el tema es promovido a sujeto. & media & proceso \\
  \bottomrule
\end{longtable}}
\pendiente{Revisar el análisis de \textit{se debe}: \textit{deberse} es más bien
cuasirreflejo inherente}
\end{solucion}

\begin{solucion}{\ref{ej:10-alternantes}}
\textit{Aclaración}: los ofrecidos en el cuadro son ejemplos posibles de uso de
los verbos.

a)

{\small
\begin{tabularx}{\linewidth}{@{}l L l L@{}}
  \toprule
  Forma transitiva & Significado & Forma pronominal & Significado \\
  \midrule
  \forma{cruzar} & \sig{pasar de un lado a otro opuesto}: \forma{Cruzamos el
    puente para llegar a Zárate.} & \forma{cruzarse con} & \sig{encontrarse con
    alguien viniendo de lugares opuestos}: \forma{Me crucé con mi vecino en el
    supermercado.} \\
  \forma{hallar} & \sig{encontrar}: \forma{No pudieron hallar la solución.} &
    \forma{hallarse en} & \sig{estar en algún lugar o situación}: \forma{En el
    sueño me hallaba en un lugar maravilloso.} \\
  \forma{aprovechar} & \sig{utilizar algo útilmente, obtener beneficio de
    algo}: \forma{Aprovechamos las vacaciones para conocer Córdoba.} &
    \forma{aprovecharse de} & \sig{sacar provecho de algo o de alguien}:
    \forma{Se aprovecharon de su ingenuidad y lo estafaron.} \\
  \forma{empeñar} & \sig{dejar en prenda algo}: \forma{Empeñó su palabra y todos
    confiaron en que cumpliría lo pactado.} & \forma{empeñarse en} &
    \sig{esforzarse, obstinarse en algo}: \forma{Está empeñado en aprender a
    nadar.} \\
  \bottomrule
\end{tabularx}}

b) Los verbos de 1 son cuasirreflejos alternantes con cambio de significado; los
verbos de 2 son alternantes, pero no presentan cambio de significado.
\end{solucion}

\begin{solucion}{\ref{ej:10-inherentes}}
En 2, 3 y 5, la introducción del pronombre produce cambios en la estructura
sintáctica: la variante sin pronombre exige un COD (\forma{sus alumnos, el
bosque, el clavo}), mientras que la variante pronominal exige un CPR. Se trata
de construcciones cuasirreflejas alternantes. Significado: en 2a, quienes se
retiran son \forma{los alumnos}; en 2b, quien se retira y saluda o se despide es
\forma{la profesora}. En 3a, \forma{cruzar} se interpreta como \sig{atravesar};
en 3b, \forma{cruzarse con} se interpreta como \sig{encontrarse con (de
casualidad), toparse}. En 5a, \forma{asegurar} se interpreta como \sig{afirmar,
reforzar, fortalecer}; en 5b, \forma{asegurarse} se interpreta como
\sig{cerciorarse, confirmar, estar seguro de algo}. En 4, la variante sin
pronombre es agramatical, lo que muestra que se trata de un caso de cuasirrefleja
pura o inherente, en la que el pronombre forma parte del lexema.
\end{solucion}

\begin{solucion}{\ref{ej:10-correr}}
En todos los casos, la introducción del pronombre produce cambios en la
estructura eventiva y afecta la delimitación del evento. 2. a) En 2b, \forma{se}
tiene valor aspectual y hace foco en el inicio del evento: se trata de una
construcción ingresiva o incoativa que produce cambio de significado: en 2a,
\forma{caer} refiere a \sig{llegar de improviso a un lugar}, pero en 2b
\forma{caerse} se interpreta como \sig{desplomarse en el suelo}. b) En ambas
oraciones, \forma{gente} tiene el papel temático de paciente. 3. a) En 3b, el
clítico no es posible porque el evento de llover se presenta como no delimitado:
se interpreta como \sig{movimiento descendente, hacia abajo} y no implica
\sig{desprendimiento de ningún lugar}. b) En 3a, \forma{la lluvia} es tema. 4. a)
En 4a no hay coincidencia entre la persona del pronombre \forma{me} y el verbo
\forma{cayó}: no hay construcción pronominal; 4b es una construcción pronominal y
tiene significado accidental. b) En 4a, \forma{me} tiene el papel temático de
meta y en 4b \forma{me} indica al afectado por el evento: es un benefactivo
negativo omisible. 5. a) En 5a no hay coincidencia entre la persona del pronombre
\forma{me} y el verbo \forma{cayó}: no hay construcción pronominal; 5b es una
construcción pronominal cuasirrefleja. b) En 5a, \forma{me} indica al afectado
por el evento, es un benefactivo negativo omisible, y \forma{la silla} es tema;
en 5b, \forma{la silla} es origen.
\end{solucion}

\begin{solucion}{\ref{ej:10-cocina}}
\forma{Ampliarse}: uso reflexivo del verbo transitivo \forma{ampliar}, que se
construye con agente y tema: \sig{alguien amplía algo}. En el texto se plantea
que H.~Milberg y su familia se amplían a sí mismos al ampliar su vivienda, que es
considerada una posesión inalienable (v. nota~\ref{fn:10-inalienable}).
\forma{Se} funciona como COD. \forma{Mudarse}: se trata de una construcción
cuasirrefleja alternante, que produce un cambio en la estructura sintáctica: el
COD se convierte en CPR y se pasa de la construcción transitiva \forma{mudar
casa} (registrada en el \textit{DLE}, aunque no utilizada en la variedad
rioplatense) a otra de transitividad indirecta, \forma{mudar(se) de casa}.
\forma{Quedarse}: \forma{se} produce un cambio en la estructura eventiva: la
variante pronominal del verbo tiene interpretación ingresiva debido al efecto
del \forma{se} delimitador.
\end{solucion}

\begin{solucion}{\ref{ej:10-intensificacion}}
2. Construcción cuasirrefleja pronominal de intensificación producida por la
incorporación de \forma{se} y el intensificador \forma{todo}. Se interpreta que
el evento se lleva a cabo hasta un punto extremo. 3. a) Construcción
cuasirrefleja pronominal culminativa de delimitación. El núcleo del COD, que es
un nombre de materia (\forma{carne}), está delimitado por el cuantificador
\forma{todo} y habilita la incorporación del \forma{se} aspectual. b)
Construcción cuasirrefleja pronominal de intensificación producida por la
incorporación de \forma{se} y el COD intensificado por \forma{hasta}.\footnote{En
estas construcciones, \forma{hasta} no es preposición, sino un focalizador, que
equivale a \forma{incluso}. Los focalizadores pueden resaltar o enfatizar
distintos constituyentes de la oración.} Se interpreta que el evento se llevó a
cabo hasta un punto extremo. 4. a) El COD de \forma{fumar} queda implícito y se
interpreta como \sig{alguna sustancia fumable, por lo general, cigarrillos}. b)
Construcción cuasirrefleja pronominal de intensificación producida por la
incorporación de \forma{se} y el añadido del COD con sentido superlativo,
\forma{la vida}. c) Construcción cuasirrefleja pronominal culminativa de
delimitación. El COD está cuantificado y permite la incorporación del
\forma{se} aspectual. 5. a) Oración intransitiva. El verbo significa \sig{ir
afuera}. b) Construcción cuasirrefleja de delimitación ingresiva. El verbo
cambia el significado a \sig{desprenderse}, con presunción de origen. El
argumento de origen es \forma{de la camisa}. c) \forma{Salirse} es una variante
pronominal de \forma{salir}. \forma{Me} es omisible y \forma{de la dieta} indica
origen (el lugar figurado de donde se sale). 6. \forma{Ir} es un verbo
intransitivo de movimiento que presupone un origen no explicitado y una meta
(\forma{a la iglesia}). Se trata de una construcción cuasirrefleja de
delimitación ingresiva.
\end{solucion}

\begin{solucion}{\ref{ej:10-diferencias}}
2. En a, el verbo transitivo tiene lectura genérica y la construcción es
impersonal; en b, la enumeración de los diferentes platos del menú se interpreta
como \sig{se comieron todo} y la construcción es cuasirrefleja de
intensificación. 3. En a, el pronombre \forma{se} produce un cambio en la
estructura argumental: se elimina el agente. Dado que el verbo está en presente y
lo acompaña el SP \forma{con facilidad}, se produce una lectura genérica y la
construcción es medio-pasiva. En b, la incorporación del pronombre produce un
cambio en la estructura eventiva: el \forma{se} delimitador convierte el evento
durativo \forma{aprender} en uno delimitado. 4. En a, el pronombre \forma{se}
provoca un cambio en la estructura argumental: se elimina el agente y se
promueve el tema (\forma{turnos de vacunación}) a sujeto de la oración. La
construcción es cuasirrefleja pasiva. En b, el pronombre \forma{se} materializa
uno de los argumentos del verbo, la meta, y cumple función de COI. Se trata de
una estructura refleja recíproca, que admite el refuerzo \forma{mutuamente /
unos a otros}.
\end{solucion}

\begin{solucion}{\ref{ej:10-akasha}}
\forma{(El aire) me rodea; (el aire) me eleva; (mi acompañante) me mira; (un
tren moderno) me despierta una inmensa curiosidad; no me importa (eso); (la
puerta) me saca del ensueño; me encanta ese lugar} y \forma{el silencio me hace
sentir incómoda}: a) la referencia del pronombre no coincide con la del sujeto y
la construcción no es pronominal.

\forma{Me pregunto (\ldots)}: a) la referencia del pronombre coincide con la del
sujeto y la construcción es pronominal. b) El pronombre manifiesta al paciente y
cumple la función de COI. i) La construcción es refleja pura porque acepta el
refuerzo \forma{a mí mismo/a}.

\forma{(El poder espiritual de carácter omnipresente) se encuentra\ldots}: a) la
referencia del pronombre coincide con la del sujeto y la construcción es
pronominal. b) El pronombre no manifiesta ningún argumento del verbo. ii) El
pronombre produce cambio en la estructura argumental (elimina al agente y
promueve al tema a sujeto de la oración con lectura de afectación). La
construcción es cuasirrefleja pasiva porque admite tercera persona singular y
plural: \forma{Los poderes espirituales de carácter omnipresente se encuentran}.

\forma{Se forma el espíritu; mis dedos se deslizan por el espacio; (mis brazos)
se convierten en alas; la puerta se abre} y \forma{(la puerta) se abrió; no me
levanto; mi sonrisa se borra; el tren apenas se mueve}: a) la referencia del
pronombre coincide con la del sujeto y se trata de construcciones pronominales.
b) El pronombre no manifiesta ningún argumento del verbo. ii) El pronombre
produce cambio en la estructura argumental (elimina al agente y promueve al tema
a sujeto de la oración, con lectura de proceso). Las construcciones son
cuasirreflejas medias porque admiten el refuerzo \forma{solo/a/os/as} o \forma{por
sí solo/a/os/as}.

\textit{Aclaraciones}: 1. \forma{Convertir} es un verbo de transformación, que
exige tres papeles temáticos ---agente, tema y meta---: \forma{Yo convierto mis
brazos en alas} (voz activa) $\rightarrow$ \forma{Mis brazos se convierten en
alas} (voz media). La meta es un CPR. 2. Con los verbos de cambio de posición o
ubicación, como \forma{levantar}, cuando el sujeto es animado se ha considerado
que es a la vez agente y paciente.
\end{solucion}
